\documentclass[lualatex,a4paper,10pt]{jlreq}
\usepackage{../../../templates/survey/jgaisurvey}
\usepackage{needspace}
\usepackage{booktabs}
\usepackage{tabularx}
\usepackage{array}
\usepackage{multicol}
\usepackage{xurl}
\addbibresource{references.bib}

\setlength{\columnsep}{1.45em}
\setlength{\multicolsep}{0.65em}

\surveysetup
  {SP-efficient-llm-2026}
  {Japanese Generative AI Technical Survey Special}
  {Efficient Intelligence――LLMを速く、軽く、安くする技術史 / Thematic survey as of 2026-09-21}
  {Open history, evidence observed through 2026-09-21 UTC}
\surveyeditiondescriptor{Thematic Longform Technical Survey}

\surveycoverstory
  {Efficient Intelligence}
  {LLMを速く、軽く、安くする技術史。訓練の配分、疎な計算、記憶の節約、低ビット、逐次復号の緩和、運用の束ね、安価計算への振り分けを、層の重なりとして読む。}
  {Bottleneck Contract \quad / \quad Conditional Compute \quad / \quad Attention \& Memory \quad / \quad Precision \quad / \quad Decoding \quad / \quad Serving \quad / \quad Routing \quad / \quad 2026 Capstones \quad / \quad Measurement Literacy}

\begin{document}
\surveycover
\clearpage

\section*{本巻の問いと読み方}
本巻の問いは、\textbf{大規模言語モデル（LLM）の効率化が、単一軸の最適化から、学習・構造・記憶・復号・運用・振り分けにまたがる層状の設計問題へと、どのように進化してきたのか}である。対象は基礎的な規模則から疎な混合専門家、注意機構と記憶保持、低ビット化と局所実行、逐次復号の高速化と推論予算の管理、サービング運用、安価計算への振り分け、そして二〇二六年の四つの実装点へと至る。本巻は製品の順位付けでも、特定モデルの推奨でも、条件の異なる数値の横断比較でもない。\textbf{どの隘路を、どの作業負荷と装置条件のもとで、何で削り、何を代償にしたか}を追う技術史である。

重要なのは、ここで「効率」を一語で済ませないことである。学習に要する計算量を減らしたのか、配備後の応答を速くしたのか、同時接続の捌き量を増やしたのか、ひとつの仕事を終えるまでの総費用を下げたのか。これらを区別せずに手法だけを並べれば、あらゆる比較は意味を失う。第1節で結ぶ用語と隘路の契約が、本巻全体の共通言語になる。

\subsection*{確認範囲と主張境界}
本巻の技術的主張は、受理された証拠の範囲に縛られる。一次資料の本文まで読み込めた箇所もあれば、要旨の範囲に限定される箇所もある。概要水準の主張は概要水準として扱い、本文で裏付けられた言語では記述しない。提供元による速度・価格・精度の数値は、測定条件付きの主張として扱い、独立検証のない数値は事実として引用しない。異なる条件で測られた性能値を直接並べて優劣を論じることはしない。測定条件や文脈長や同時実行数が異なれば、律速要因そのものが変わるからである。

未解決の事項は未解決のまま残す。独立した再現が未了の主張、公開の較正手順が見当たらない評価、単独の製品姿が確認できない後継の存在、本文未読の制約、版や対応状況の未確定は、記事を読みやすくするために閉じない。境界の明示は弱さではなく、本巻の読みの作法である。

\subsection*{この巻で比較するもの}
本巻が横断的に比較するのは、モデルの絶対順位ではなく、\textbf{どの設計問題が第一級の競争軸になったか}である。具体的には、学習の資源配分、疎な条件付き計算、注意と記憶保持の設計、数値精度と局所実行、逐次復号と推論予算、運用の束ねと受け渡し、安価計算への振り分けと特化判断を追う。二〇二六年の四つの実装点は、順位付けではなく機構の積み重ねとして読み解く。最終節では、数字を読むための物差しそのものを点検する。

\subsection*{読者への道案内}
時間の限られた読者は、第1節の契約表と各節末の境界ボックスだけを追っても、本巻の主張の骨格をつかめる。各節の用語は第1節の区別表に準拠する。総量と有効量、重みとキャッシュ、トークン単価と仕事単価の三つの分離が、各節で繰り返し使われる合言葉である。実装の詳細よりも設計思想の違いに関心のある読者は、第2節の系譜表、第3節の節約位置表、第8節の比較表を中心に読むとよい。運用と測定の実務に関心のある読者は、第6節の無駄と処方の表、第7節の設計表、第9節の作法表が近道になる。巻末の技術要点表は、各節の主張と境界を一次資料の裏付けとともに圧縮したものであり、本文の詳述に戻る際の索引として使える。巻末の用語集は、本巻で使う技術用語の定義を集めたものであり、節をまたいで同じ言葉が同じ意味で使われていることを確かめるために置いた。

\setcounter{tocdepth}{2}
\tableofcontents
\clearpage

\section{「効率」とは何を減らすことなのか}
\label{sec:bottleneck}
\sectionkicker{BOTTLENECK CONTRACT}

\begin{multicols}{2}
\raggedcolumns
効率が良いという言葉は、大規模言語モデルの文脈では、それだけでは何も約束しない。学習に要する計算量を減らしたのか、配備後の推論の応答を速くしたのか、同じ費用でさばける要求数を増やしたのか。これらを区別せずに手法だけを並べれば、後のあらゆる比較は意味を失う。本稿が最初に結ぶのは、その混乱を避けるための契約、すなわちどこが詰まるのかを見定める見取り図である。\autocite{kaplan2020,chinchilla,orca}

出発点は学習と推論の分離である。学習は一度きりの巨大な投資であり、モデル規模と学習データ量と投下する計算量の配分が問われる。損失をどこまで下げるかが中心であり、使った資源は一度支払えば済む。これに対して推論は、配備後に繰り返し発生し続ける継続費用であり、使われ方によって効き方が変わる。同じ速さであっても、ひとつの応答が返るまでの待ち時間である遅延と、単位時間あたりに処理できる要求量であるスループットは別の軸である。対話の応答性を高めることと、同時接続の捌き量を増やすことは、同じ工夫では解けない場合が多い。さらに、トークンあたりの単価を下げたことと、ひとつの仕事を終えるまでの総費用を下げたことも等しくない。短い生成を安く大量に回せたとしても、長い文脈の読解や複数段階の推論を要する仕事の費用が下がらなければ、利用者にとっての効率は改善しない。

計算か記憶容量か帯域か通信か、総パラメータか有効に働く部分か、固定資産としての重みか文脈に応じて膨らむ生成時の中間記憶か。総量の大きさと実際に働く有効部分の大きさを混同せず、固定資産と変動費を分けて考える習慣が、以降の節約策を見極める土台になる。この契約なくして高速化や軽量化を語れば、遅延の改善か捌き量の改善か、固定費の削減か変動費の削減かが曖昧になり、読み手は判断を誤る。削減対象が違えば選ぶべき節約策も評価の物差しも変わるという当然の事実から出発したい。何を減らすのかを名指しすることこそ、本特集の効率論の出発点であり、以降の手法を読み解く際の共通言語なのである。

\subsection{学習の物差しとしての規模則とその見直し}
学習側の物差しとして長く参照されてきたのは、言語モデルの損失がモデル規模とデータ量と計算量に対してべき乗則に従って滑らかに減少するという見方である\autocite{kaplan2020}。ここで重要なのは係数の具体値ではなく、大きくすれば損失が下がるという関係の形そのものである。この形が、濃密なモデルを前提とした規模拡大の議論の土台になり、研究開発の資源配分を正当化してきた。まず大きなモデルを作り、後から小さく使うという発想も、この方向性の共有に支えられている。しかしこの関係は、あくまで学習時の損失に関する見取り図であり、配備後の推論費用の内訳を解消するものではない。規模だけを信奉すれば、保持すべき重みは増え、それを読み出す帯域の負担も増し、複数装置にまたがる通信の重さも増す。演算器の処理能力と、記憶容量と、読み出しの帯域と、装置間の通信は、それぞれ異なる局面で律速になる。損失の低下と引き換えに配備の重さが増すという緊張関係がここに生まれる。本稿が依拠する範囲では、係数の具体値や実験条件の細部には立ち入らず、べき乗則という関係の形を取り出すにとどめる。

学習の資源配分を問い直したのは、従来の大型モデルは投じた規模に見合う学習トークンを使い切れておらず、過小学習の状態にあったという指摘である\autocite{chinchilla}。器だけを大きくし、中身の学習が追いついていない状態では、投下した計算量が見合いの損失低下に結びつかない。そこでモデル規模と学習トークン数をほぼ均等に伸ばす配分が代表的な処方箋として示され、パラメータあたり約二〇トークンという目安が広く知られるようになった。肝心なのは、この数値を普遍の定数として扱わないことである。本稿が依拠する範囲では、それは特定の条件のもとで導かれた見出しとしての提案であり、条件や構成が変われば望ましい配分も変わり得るという留保つきで受け止める。むしろ要点は、効率の問いがモデルの大きさだけでなく、データの割り当て方へと移った点にある。同じ計算量を使うなら、巨大な器を不十分なデータで鍛えるよりも、規模とデータの釣り合いを取り直す方が損失の低減に資するという視点転換である。総量と有効量を区別し、トークンあたりの単価ではなく、ひとつの仕事を終えるまでの費用で考える習慣がここから要請される。

\subsection{推論の詰まりを分けて見る}
推論側の詰まりを可視化したのは、生成過程を一括処理から反復単位の細粒度な管理へと変えた発想である\autocite{orca}。生成は、入力全体を一気に読み込んで内部状態を構築する前処理であるプリフィル段階と、トークンをひとつずつ逐次的に生み出す復号であるデコード段階に分かれる。前者は並列に処理しやすく演算が律速になりやすいのに対し、後者は直前の結果に依存して逐次に進むため、重みの読み出しや蓄積された中間記憶の保持と読み出しが律速になりやすい。従来のように要求をひとまとまりとして最後まで占有させる方式では、長い生成が短い要求を待たせ、遅延とスループットの両方が損なわれる。そこで導入されたのが、反復のたびに待ち行列を見直し、段階の異なる要求をきめ細かく詰め合わせる継続的なスケジューリングの考え方である。分散した装置群の上でも反復単位の管理を徹底する発想は、その後のスケジューラ設計の起点になる考え方であり、段階の分離が運用の柔軟性を生むことを示した。

ここで効く区別が、重みと中間記憶の違いである。重みはモデルに固有の固定資産であり記憶容量を占有する。一方、逐次生成に伴って蓄積されるキーとバリューの中間記憶であるKVキャッシュは、文脈の長さや同時処理数に応じて膨らむ変動費であり、容量と帯域の両面から圧迫する。計算能力と記憶容量と帯域と通信のどれが詰まっているのかを名指しし、遅延とスループットのどちらを改善したいのかを明確にすること。それが推論の効率化の前提であり、手法の評価軸を定める作業にほかならない。本稿では、反復単位でプリフィルとデコードを分けて扱うという設計思想の水準にとどめ、具体的な処理手順の細部や数値化された高速化の幅には踏み込まない。

\subsection{総量と有効量の分離が効く場面}
総パラメータ量と有効パラメータ量の区別は、学習の議論と推論の議論を切り分けるための基本的な物差しである。総パラメータ量はモデルが保持する重みの全体規模を表し、保存に必要な容量や読み込みの対象となる量を決める。一方、有効パラメータ量は一つのトークンを処理する際に実際に計算へ寄与する部分の規模を表し、一トークンあたりの計算量の見通しを与える。この二つを分けないまま「大きいモデル」「小さいモデル」と呼ぶと、保存や転送の重さの話と計算の重さの話が混ざり、どの対策が効くのかが見えなくなる。総量を問題にするのか有効量を問題にするのかを先に宣言すること自体が、本節の用語の使い方の核心である。

この分離がとくに効くのは、重みの大きさとキャッシュの大きさを別々に考えたい場面である。重みは入力によらず固定された量として存在し、モデルの総量によって決まる。これに対してキーとバリューのキャッシュは、処理する文脈の長さや同時に抱える要求の数に応じて増減する。短い応答を一つだけ返す場合には重みの読み出しが目立ち、長い文脈を保ったまま多くの要求を並行してさばく場合にはキャッシュの置き場と読み書きが目立つようになる。同じモデルを使っていても、利用の形が変われば重みとキャッシュのどちらが先に苦しくなるかは変わる。だからこそ、重みの量とキャッシュの量を一つの「メモリ使用量」という言葉に丸めず、総量に由来する分と利用条件に由来する分に分けて呼ぶことが重要になる。

トークンあたりの費用と課題あたりの費用の区別も同じ発想である。一トークンを処理する費用が下がっても、課題を解くために必要なトークンが増えれば、課題全体の費用は下がらない。たとえば中間的な推論過程を長く書き出す手法や、複数回の試行から良いものを選ぶ手法では、一回ごとの生成は軽く見えても、課題を解き終えるまでに消費するトークンの総量は増える。逆に、一トークンあたりの費用が多少上がっても、必要なトークンが十分に減れば課題全体では安くなる場合がある。効率の主張を読むときは、分母がトークンなのか課題なのかを必ず確かめる必要がある。分母の取り違えは、軽量化の効果の過大評価にも過小評価にも直結する。

規模則の議論は、この分母の問題を学習側から照らす。言語モデルの損失がモデル規模やデータ量や計算量に対して滑らかなべき乗則に従うという見方は、闇雲な拡大ではなく配分を考える土台を与えた\autocite{kaplan2020}。その後の見直しは、過大なモデルに従来より多くの学習トークンを割り当てる配分が学習効率の観点で有利になり得ることを示し、効率の問題をモデルの大きさの問題からデータ配分の問題へと広げた\autocite{chinchilla}。ここでいう有利さは学習計算の配分に関する主張であり、推論時の重みの重さやキャッシュの重さまで含めた総合判断ではない。学習の節で使う「効率」と推論の節で使う「効率」の分母が違うことを、読者はここで予告として受け取ってほしい。総量と有効量、重みとキャッシュ、トークンと課題という三つの分離は、その予告を具体的な言葉に直したものである。

\subsection{測定条件が変われば律速が変わる}
同じモデルと同じ実装を使っても、測定条件が変われば律速は変わる。律速とは、全体の速さや費用を決めている支配的な要因のことであり、計算の速さ、保持できる量、読み書きの速さ、並列実行に伴う通信の四つに大別できる。効率化の手法は、この四つのうちどれを緩和するかが違えば、効く条件も違ってくる。手法の名前だけを見て優劣を決めるのではなく、どの律速を緩める手法なのかを問う習慣が、本特集を通じた読み方の約束になる。

最も分かりやすい条件の違いは、一度に処理する量である。一つの要求への応答の速さを問題にする場合、待ち時間の短さが評価の中心になる。利用者は最初の文字が出るまでの速さと、その後の続きが出る滑らかさを感じる。この条件では、装置を十分に使えていない状態での無駄や、処理の切れ目での待ちが律速になりやすい。一方、多くの要求をさばく処理能力を問題にする場合、単位時間に完了する要求の数が評価の中心になる。この条件では、複数の要求を束ねて装置の稼働率を上げる工夫が効き、束ねた分だけキャッシュの置き場や読み書きの負担が増える。待ち時間の話と処理能力の話は、同じ装置の同じモデルでも結論が逆になることがあるため、分けて記録しなければならない。

文脈の長さも律速を動かす。入力を一気に読み取る段階では、まとまった計算を一括して実行できるため計算の能力が問われやすい。これに対して出力を一語ずつ生成する段階では、その都度重みとこれまでのキャッシュを読み出す必要があり、読み書きの帯域が問われやすい。文脈が短いうちは重みの読み出しが相対的に目立ち、文脈が長くなるにつれてキャッシュの読み書きの比重が高まる。この段階の違いに着目し、入力読み取りと逐次生成のスケジューリングを分けたことが、生成モデル向けサービングの出発点にある\autocite{orca}。前処理と生成を同じ待ち行列で一括して扱うか、段階ごとに分けて扱うかという設計の違いは、待ち時間と処理能力のどちらを優先するかに直結する。

同時実行の数と装置構成も同様に律速を動かす。同時に抱える要求が増えれば、キャッシュの総量は膨らみ、保持できる量が先に上限になることがある。複数の装置に分けると、一台あたりの保持量は緩和されるが、装置間の通信が新たな律速として現れる。帯域の広い装置に替えれば読み書きの律速は緩むが、計算の能力が追いつかなければ今度は計算が律速になる。つまり律速はモデルの属性ではなく、模型と利用条件と装置の組み合わせの属性である。効率の比較は、束ねる量や同時実行数や文脈長や装置の別をそろえたうえでなければ意味を持たない。条件を明示しない比較は、本特集では成立しない比較として扱う。

後の学習効率の議論では、総パラメータ量と学習トークンの配分、計算量の使い方といった言葉が中心になる。損失の縮小を目的とする学習の場面では、重み全体の規模とデータ配分の関係が問われ、推論時のキャッシュや同時実行の話は背景に退く。これに対して後の推論基盤の議論では、有効パラメータ量、重みとキャッシュの分離、入力読み取りと逐次生成の区別、待ち時間と処理能力の区別が中心になる。軽量化や量子化の議論では、重みの保持量と読み書きの帯域、計算の簡略化の関係が中心になる。同じ「効率」という語が節ごとに違う分母と違う律速を指すことを、読者は各節の冒頭で確かめながら読み進めてほしい。

\end{multicols}

\begin{center}
\small
\begin{tabularx}{\textwidth}{@{}p{0.18\textwidth}p{0.36\textwidth}X@{}}
\toprule
区別の軸 & 両側の意味 & 本巻での使い方 \\
\midrule
学習と推論 & 一度きりの投資と繰り返しの継続費用 & 規模則・配分の議論と運用の議論を混同しない \\
プリフィルとデコード & 並列に読む段階と逐次に生む段階 & 演算律速と記憶帯域律速を分けて読む \\
計算と容量・帯域・通信 & 演算器・置き場所・読み出し・受け渡し & 律速の所在を名指しする \\
遅延とスループット & 応答の速さと捌き量 & どちらを改善したのかを明示する \\
総量と有効量 & 総パラメータと稼働パラメータ & 大きさの印象を稼働の実態で割り引く \\
重みとKVキャッシュ & 固定資産と変動費 & 容量見積もりを分けて立てる \\
トークン単価と仕事単価 & 単位の安さと完遂の費用 & 費用の物差しを完遂側に置く \\
\bottomrule
\end{tabularx}
\end{center}
\autocite{kaplan2020,chinchilla,orca}

\begin{multicols}{2}
\raggedcolumns
以上の見取り図は、あえて限定された範囲の主張だけに依拠している。規模と損失の関係については関係の形を取り出すにとどめ、学習トークンの配分については代表的な目安を見出しとして紹介するにとどめ、生成の管理方式については設計思想の水準にとどめる。異なる条件で測られた性能値を直接並べて優劣を論じることも、本稿では行わない。本稿が約束するのは削減対象の名指しである。以降の各論は、この物差しに沿って読まれることを前提とする。適用範囲を超えた一般化を避け、条件付きの主張は条件付きのまま受け止めること。それが本特集の効率論を読み解く際の作法であり、本稿の守備範囲の宣言である。
\end{multicols}

\begin{claimboundary}[本節の読解上の境界]
規模則の係数値・実験条件、配分目安の普遍化、処理手順の細部と速度幅には立ち入らない。異なる土俵の数値を並べた速さ比べは行わない。学習の物差しを推論の評価に流用しない。
\end{claimboundary}

\section{ScalingからConditional Computeへ}
\label{sec:moe}
\sectionkicker{CONDITIONAL COMPUTE}

\begin{multicols}{2}
\raggedcolumns
言語モデルの規模拡大は、パラメータとデータと計算をそろえて増やすほど性能が伸びるという経験則に支えられてきた。緻密なモデルではすべてのトークンがすべてのパラメータの計算を負担するため、能力を上げようとすればそのまま計算量と費用が膨らみ、記憶容量や装置間通信の壁にも早くぶつかる。この行き詰まりをほどいたのが、容量と計算を切り離す条件付き計算の発想である。学習されたゲーティングがトークンごとにごく一部の専門家だけを選び出し、残りは休ませることで、総パラメータ数を増やしてもトークンあたりの計算を抑えられる\autocite{shazeer2017}。均等に割り振るための補助損失が添えられ、特定の専門家への偏りを抑える仕組みも一体で示されたことが重要である。ここでは細かな数式や長期系列での実験表には立ち入らず、仕組みの要点だけを調査水準で受け取る。

この着想を巨大規模で現実にする方向へ進めたのが、専門家並列と自動分割の系統である。専門家を複数装置に分散配置し、トークンを振り分けて処理することで、単一装置の記憶容量を超える規模が扱えるようになった\autocite{gshard}。振り分け先があふれた場合の容量係数による扱いや、分割をコンパイラが管理可能にする工夫が、後の兆級モデルの学習基盤の前提となった。翻訳ベンチマークの細部や分割注記の全体には踏み込まず、疎な拡大を可能にした体制側の前進として押さえる。総量と稼働量を分けて考える習慣もここに始まり、後の稼働量と常駐量の経済学へとつながる。同時代の実装点同士の優劣比較はここでは行わず、系譜と割当の仕組みに集中する。

\subsection{振り分けの設計問題}
疎な層をどう振り分けるかは、質と通信と運用の三つ巴の設計問題になった。振り分けを一人だけに絞る方式は、仕組みを単純化し、装置間のやり取りと割り当ての管理を軽くする道を開いた\autocite{switcht}。トークンあたり一人か複数かという選択は、以降の細分化や共有専門家を評価する際の基準となり、方式の系譜を整理する物差しとして残った。ここでは規模拡大曲線や多言語の評価表には踏み込まず、単純化が何を節約し何を手放したかという交換関係だけを調査水準で受け取る。対照的に、専門家側がトークンを選ぶ方式は、視点を反転させたもう一つの到達点である\autocite{expertchoice}。各専門家が受け持てる量を固定し、トークンごとに選ばれる数が変わるため、容量の管理が明確になり、負荷の偏りへの向き合い方が変わる。事前学習の高速化の定量分析には立ち入らず、方式の違いが通信量や待避や常駐のどこに効くかという水準で整理する。いずれも万能ではなく、装置構成や系列長や運用条件で損得が変わる点に注意が要る。

\subsection{細分化と共有化、そして稼働量の経済学}
専門家を細かく砕き、共通の知識を担う共有専門家を添える設計は、条件付き計算の使い方を一段と経済的にした\autocite{deepseekmoe}。細分化された専門家が専門語彙のように振る舞い、汎用的な部分は共有部が引き受けることで、組み合わせの表現力と通信の節度を両立させる狙いがある。装置をまたぐ転送を抑えるため、振り分け先の装置数を絞る工夫も一体で示された。ここでは専門家粒度の走査などの詳細な比較表には踏み込まず、抽象水準の仕組みとして受け取る。この機構を受け継いだ分岐点として、注意機構の圧縮と疎な計算を同時に前面に出した大規模モデルが現れた\autocite{deepseekv2}。総量に対して稼働量を抑える姿で示され、鍵と値の保持を圧縮する機構と疎な演算を組み合わせることで、記憶と計算の双方に手を打つ姿勢が明確になった。定式化の全体や評価表の数値には立ち入らず、確立した水準の構成的事実として押さえる。

総量と稼働量を分けて見る習慣が定着すると、巨大さの印象は大きく変わる。開かれた重みで広く使われた方式は、8x7B構成でトークンごとに二つの専門家を選び、稼働量を抑える姿として確立した水準で知られた\autocite{mixtral}。この公開を境に、専門家の常駐や待避、振り分けの上乗せ費用といった運用の論点が共有の課題となり、巨大な重みを手元でどう保持し、必要な分だけをどう載せ替えるかという観点が効率の議論に定着した。計算の節約は記憶の負担と引き換えであり、動かさない重みもどこかに置かねばならず、通信と保持の費用が残る。単純な大小比較は意味を持たない。

\subsection{学習を支える土台}
疎な大規模化が現実の学習として回るには、低精度演算の土台が欠かせない。総量671Bの混合専門家モデルで、混合精度による大規模学習が実生産として示されたことは、その結節点である\autocite{deepseekv3}。行列演算庫の側では、細粒度の低精度演算を担う実装がこの学習の片翼となり、形式の工夫が大規模学習の実行可能性に直結することが示された\autocite{deepgemm}。同じ系譜では、複数トークンの予測を学習目的に取り込む試みも示され、後の推論高速化の起点として位置づけられる。あくまで目的関数の由来として押さえ、速度の定量主張には踏み込まない。

巨大モデルを学習可能にする分割と配置の工夫は、計算を減らす技術ではなく、載せて動かすための技術として区別して理解したい。層内を複数装置に割り付ける並列化は、より大きな模型を収めるための方向であり\autocite{megatron}、最適化状態や勾配や重みの重複を排して学習可能な規模を伸ばす記憶最適化も、収めるための方向として整理される\autocite{zero}。本特集では配置と割当の文脈に限定して扱い、分散学習の深掘りにはしない。最適化手法の側では、勾配更新に直交化を用いる実装が参照対象となり、学習の安定や高い学習率の許容が報告として伝えられた\autocite{muon}。ただし参照先は論文ではなく実装であり、収束定理の根拠として扱わない。

量を積む方向と並んで、何をどれだけ混ぜるかという構成の方向が、近年の学習効率のもう一つの柱になった。小規模な代理模型で領域ごとの重みを定め、下流課題の知識なしに混合比を決める試みは、量の拡大を補う構成の工夫として位置づけられる\autocite{doremi}。知識の蒸留は、教師の柔らかな出力を手がかりに小規模模型へ能力を受け渡す古典的な出発点であり\autocite{hinton2015}、低ランク適応は事前学習済みの重みを凍結し低ランクの行列だけを学ぶことで微調整の記憶と保存を軽くする手法として知られる\autocite{lora}。いずれも本特集では深掘りではなく、系譜の指し示しとして配置する。効率を一枚岩で論じず、局面ごとの費用として分けて考える習慣がここで整い、次の記憶と復号の議論へとつながる。

\subsection{学習の割当としてのデータ組成}
規模則の議論は計算量とデータ量の量的配分に目を向けがちであるが、何を混ぜて学ばせるかという組成の設計もまた割当の問題である。混合比の設計は、トークン選択型の疎なゲーティングが計算と容量を切り離したのとちょうど対称的に\autocite{shazeer2017}、データ側の計算と容量の切り離しを担う。すなわち、同じトークン予算でも組成が違えば到達する能力分布が変わり、疎なモデルにおける専門家の分化のさせ方も変わる。この意味でデータ組成は前処理の職人芸ではなく、条件付き計算と並ぶもう一つの割当機構である。

この割当を原理的に扱おうとした試みが、領域別の混合重みを事前に設計する考え方である。小規模な参照模型を用いて領域上の頑健な目的を最適化し、下流課題の知識なしに混合比を定めてから大規模な本学習に移す手順は、量の拡大を補う組成の拡大と理解できる\autocite{doremi}。どの領域に重みを置くかは、どの専門家を育てるかという問いと間接的に連動する。共通知識を担う共有専門家と細分化された専門家の住み分けが、粗い混合比の選択によって歪むこともあれば、整えられることもあるからである\autocite{deepseekmoe}。組成設計はしたがって、振り分けの経済学の上流にある割当問題として読むべきである。

蒸留は、この組成の議論を容量と品質の交換関係へと接続する。教師の柔らかな出力を目標として学ぶ枠組みは、硬い正解ラベルだけでは伝わらない分布の形状を移すことで、小さな容量でも高い品質に到達する道を開いた\autocite{hinton2015}。疎な大規模モデルが総パラメータと作動パラメータを切り離すのに対し、蒸留は学習済みの容量と転移可能な品質を切り離す。開かれた重みとして公開された疎モデルが下流の常駐や振り分けの関心を呼び起こしたように\autocite{mixtral}、蒸留された小さな模型は配置と推論の関心を引き寄せる。いずれも、大きな容量をどこに置き、どの部分だけを作動させるかという同じ経済の別表現である。

低ランク適応は、この系列の末端を微調整の効率として締めくくる。事前学習済みの重みを凍結し、学習可能な低ランク分解行列だけを挿入する方式は、微調整の記憶量と保存量を切り詰める\autocite{lora}。ここで節約されているのは事前学習の計算ではなく、適応のための常駐と更新の対象範囲である。条件付き計算がトークンごとの作動範囲を絞るのに対し、低ランク適応は課題ごとの更新範囲を絞る。事前学習における混合比の設計と、適応における更新範囲の設計が、予算配分の上下流として対応していることが見えてくる。

\subsection{載せるための技術と削る技術の区別}
大規模化の技術はしばしば一括りに節約技術と呼ばれがちだが、載せるための技術と削るための技術は区別しなければならない。載せるための技術は、より大きなモデルを学習可能にするための配置と並列化の前提であり、それ自体は演算量を減らさない。疎モデルの系譜でいえば、自動分割による専門家並列と振り分けの仕組みが巨大化の前提となり\autocite{gshard}、簡素化された経路選択が通信と振り分けを軽くし\autocite{switcht}、細粒度化と共有専門家が通信制約のもとでの専門化を進めた\autocite{deepseekmoe}。これらはいずれも配置と作動の設計であり、削減そのものではない。

層内並列は、載せるための技術の典型である。層内部のテンソルを行列並列に分割し、最小限の通信操作で結合する方式は、単一装置の記憶容量を超える規模の学習を可能にする\autocite{megatron}。並列化が進んでも学習の総演算は減らず、むしろ通信という新たな費用が表面化する。疎モデルではこの通信費用が専門家の振り分けと常駐の問題として現れる。載せることは削ることの対立物ではなく、削るべき対象を明確化する前提である。

冗長性の排除は、載せるためのもう一つの柱であり、並列化とは異なる方向から記憶の問題を解く。最適化器の状態や勾配やパラメータが装置間で重複保持されがちであり、この重複が規模の上限を決める。重複を排除するためにそれらを分割配置し、装置数に応じて学習可能な規模を伸ばす方式は、適合のための技術として位置づけられる\autocite{zero}。ここでも演算そのものは減っていない。減っているのは重複であり、増えているのは利用可能な総記憶である。冗長性排除は総パラメータを置く場所を確保する技術であり、条件付き計算は置いた総量のうち作動させる部分を絞る技術であることがわかる。前者がなければ後者の経済は成立しない。

直交化更新の考え方は、これら二つとは層が異なる。勾配の更新に行列の直交化操作を適用し、学習の安定性を高めてより高い学習率を許すとされる方式は、配置ではなく更新則の設計に属する\autocite{muon}。参照実装として共有された実装体に権威の所在があり、収束定理の供給源として読むべきものではない。更新則が安定化しても、専門家の常駐量や振り分け通信の総量が消えるわけではない。むしろ安定した学習が可能になることで、より細粒度な専門化への挑戦が可能になり、配置設計への要求がいっそう明確になる。

この三つを並べると、通信と常駐の見通しが自然に開ける。載せる技術が整うほど、残る律速は演算ではなく移動と滞在になる。専門家並列における振り分けの往復、容量超過時の扱い、装置制約のもとでの経路選択はいずれも移動の設計であり\autocite{gshard}\autocite{deepseekmoe}、総量に対する作動量の比率や圧縮された注意機構との組合せは滞在の設計である\autocite{deepseekv2}。微細化された量子化による行列積の土台が大規模な低精度学習を支えるように\autocite{deepseekv3}\autocite{deepgemm}、通信と常駐の土台がなければ条件付き計算の節約分は目減りする。節約は削った量ではなく、載せたうえで作動させずに済ませた量として初めて計測できる。次節以降の推論効率の議論は、この計測の作法を引き継ぐ。

\end{multicols}

\begin{center}
\small
\begin{tabularx}{\textwidth}{@{}p{0.20\textwidth}p{0.30\textwidth}X@{}}
\toprule
対象と観点 & 問い & 処方と限界 \\
\midrule
割当としての組成 & 有限予算をどの分布に振り向けるか & 参照模型による混合重みの事前設計\autocite{doremi}。最適性は分布依存であり、定式化の細部は範囲外 \\
蒸留と低ランク適応 & 容量を増やさず品質と適応を得るか & 柔らかな教師目標による転移\autocite{hinton2015}、凍結重みへの低ランク挿入\autocite{lora}。微調整効率は事前学習効率と別物 \\
層内並列による配置 & 単一装置に載らない層をどう載せるか & 層内テンソル分割と最小限通信\autocite{megatron}。総演算は不変であり、疎化の節約とは別勘定 \\
冗長性排除と直交化更新 & 重複記憶と不安定更新をどう抑えるか & 状態・勾配・パラメータの分割配置\autocite{zero}、直交化による更新整形\autocite{muon}。配置は前提であり、実装知は定理の代替にならない \\
\bottomrule
\end{tabularx}
\end{center}

\begin{center}
\small
\begin{tabularx}{\textwidth}{@{}p{0.20\textwidth}p{0.34\textwidth}X@{}}
\toprule
系譜の節目 & 設計の要点 & 残る交換関係 \\
\midrule
疎ゲート混合専門家 & トークンごとに一部の専門家だけを稼働 & 偏りと補助損失の運用 \\
専門家並列・自動分割 & 分散配置で単一装置の容量を超える & 容量係数と通信の重さ \\
単純化された振り分け & 一人選択で実装と拡張を見通す & 豊かさと見通しの交換 \\
専門家側選択 & 受け持ち量の固定で容量管理を明確化 & あふれの扱い \\
細分化と共有専門家 & 組み合わせの表現力と転送の節度 & 管理の複雑さ \\
開かれた疎構成 & 常駐・待避・載せ替えの共有課題 & 重みの置き場所 \\
低精度学習基盤 & 混合精度で大規模学習を成立させる & 数値の安定と速度 \\
分割・配置・記憶最適化 & 載せて動かすための方向 & 計算削減との混同に注意 \\
組成の設計 & 混合比・蒸留・低ランク適応 & 割当と組成の二軸 \\
\bottomrule
\end{tabularx}
\end{center}
\autocite{shazeer2017,gshard,switcht,deepseekmoe,deepseekv2,mixtral,deepseekv3}

\begin{claimboundary}[本節の読解上の境界]
同時代の実装点同士の優劣比較は行わない。数式の細部、評価表の数値、通信量の定量分析、収束定理の根拠には立ち入らない。分散学習の深掘りは範囲外であり、割当の前提としてのみ扱う。
\end{claimboundary}

\section{AttentionとMemoryを減らす}
\label{sec:attention}
\sectionkicker{ATTENTION AND MEMORY}

\begin{multicols}{2}
\raggedcolumns
長い文脈を扱うほど推論時の記憶は膨らむ。その最初の絞り込みは、鍵と値の持ち方を変えることから始まった。すべての照会頭が独自の鍵と値を抱える方式では、系列が延びるにつれて保持すべき状態が増え続け、復号の足かせになる。そこで単一の鍵・値頭をすべての照会頭で共有する方式が提案され、保持量を仕組みとして削る最も素朴な一歩になった\autocite{mqa}。品質面の代償の詳細な数値表には立ち入らず、確立した機構の水準で押さえる。続いて登場した束ねる共有方式は、複数の照会頭を少数の鍵・値頭に寄せることで、削減幅と表現力の折り合いをつけ直したものとされる\autocite{gqa}。学習済みの重みを起点に作り替えに近い形で移行できる点が実用上の鍵であり、全面学習なしに導入できる移行可能性が普及を後押しした。

一方、近似を一切入れずに速くする道もある。注意計算そのものは厳密に保ったまま、行列の分割と正規化の逐次化によって高帯域記憶との往復を減らし、実時間の効率を高める手法がそれだ\autocite{flashattn}。複雑さの証明や速度表の水準には踏み込まず、厳密さと入出力効率の両立という役割で押さえる。\textbf{計算結果を変えずに入出力の動かし方だけを磨くのか、それとも計算や保存の対象そのものを変えるのか}。この区別が本節全体の物差しになる。以降の疎化・線形化・状態削減・再利用のいずれも、どちらの側に立つのかを意識して読み分けたい。

\subsection{疎化と局所化}
すべてのトークンを見る代わりに、見るべき場所だけを選ぶ。階層的な圧縮と選択と局所走査を組み合わせ、実行環境に寄せた実装とともに最初から学習する方式は、推論時だけの刈り込みとは一線を画すとされる\autocite{nsa}。索引器の作りや共有の方式や学習形態の違いを取り違えないことが重要であり、似た名称の疎注意群の混同を防ぐことが特集全体の用語の要になる。対照的に、窓の中だけを見る局所注意は実績ある枯れた形であり、共有化された鍵・値の仕組みと組み合わせた運用例が生産的な基準になった\autocite{mistral7b}。以降の疎な大域設計や窓の再生をうたう仕組みは、この基準との対比で理解するのが見通しがよい。

新しい疎注意の登場を告げる公式発表は、告知の範囲として受け止める\autocite{dsaexp}。既存基盤の上に新しい機構を載せ、長い文脈の学習と推論の効率化をうたう内容であり、利用面の提供開始と重み・報告書・実装の公開が合わせて告げられた。索引の数式や等価性の検証手順は報告書側に属する\autocite{v32report}。発表時点の価格改定の扱いも告知の水準にとどめ、効率と価格の伝達の速さを断定しない。発表と報告を対にして読み、役割の違いを保ちたい。

\subsection{回帰と線形注意の系譜}
疎化の中身を報告書の該当節の範囲でたどる。新しい疎注意は、少数の軽量な索引機構で関連の強い受け手を選び出し、細かい粒度で注目対象を絞る設計とされる。等価性に関する報告は抜粋の範囲にとどめ、互角性の主張は変種や評価条件に縛られるものとして、確立した水準を超えて引用しない。一方、回帰的な記憶の側にも系譜がある。差分規則による記憶更新に調節の門を加えた方式は、従来の回帰型の記憶制御を強めたものとして位置づけられ、後に続く線形注意系の直接の親になると整理されている\autocite{gdn}。その子にあたる線形注意の提案は、門をさらに細かくし、区分的な計算法と全域注意との交互配置で大域の流れを保つ設計を報告する\autocite{kimilinear}。比較表の数値は該当節の範囲として扱い、独立した再現は未確定とする。報告書の利用条件が再利用に制約を課す点にも留意する。

回帰と疎注意を組み合わせた実装点も現れている。一方の陣営の公開情報は、層ごとの外部記憶表をホスト側に置き、回帰層と疎注意層を一定比率で積む構成として、要約と傍証の範囲で整理する\autocite{qwenreadme}。自己報告の評価値は検証待ちの提供者主張として隔離し、引用しない。ここで系譜の扱いが要になる。外部表を持つ設計と後に述べる条件付き記憶の提案との間に、導出関係を示す記述は収集範囲では見つかっていない\autocite{qwenlineage}。似た圧力への収束的な応答、すなわち並行ないし未確定の関係として扱い、影響関係を主張しないのが正しい。一方、線形注意側の実行可能実装は、報告を超えた配置可能性の傍証になる\autocite{kda}。ただし数値や精度の水準、性能の事実には踏み込まず、引用しない。呼称の衝突が未解決の別ラベルには触れず、混同を避ける。

\subsection{保持を削る・使い回す・深さを絞る}
保持する鍵と値を削る方向は、選ぶ仕組みを変えずに捨て方を工夫する。重要度の高い少数の受け手と直近の履歴だけを残す方式は、注意の値に効くのはごく一部のトークンだという観察に立つ\autocite{h2o}。冒頭のトークンが吸い寄せる現象への対処も欠かせない。最初の受け手の鍵と値を残し、窓と組み合わせることで窓注意の性能を取り戻す工夫が示された\autocite{sinks}。層ごとの配分を変える発想もある。層によって注意の散らばりが違うことを利用し、保持枠を層間で山形に配る設計である\autocite{pyramidkv}。本文の具体的な算法や評価の全文は未収集として、数値の比較には踏み込まない。いずれも計算式を近似するのではなく残す量と場所を絞る工夫であり、保存側の節約にあたる。

作った鍵と値を使い回すことも大きな節約になる。接頭辞の再利用を自動化する仕組みは、鍵と値の塊に印をつけて要求間で共有し、塊粒度の再利用を担う\autocite{vllmprefix}。文書の本文を直接確認していないため、確立した役割の水準で述べる。一方、特定の実行基盤から切り離して再利用や退避や分離を担う層もある\autocite{lmcache}。注意の数学を変えずに再計算を省く点で、提供時の節約にあたり、構造側の変更とは区別したい。条件付きの記憶という別の方向もある。古い埋め込み手法を現代化し、字句の圧縮や文脈に応じた調節で定数時間の参照を目指す提案である\autocite{engram}。配分曲線の詳細や学習費用の精算には踏み込まず、報告の該当範囲で整理する。実装の展示の水準には踏み込まず、紙の主張の範囲で述べる。使い回しと条件付き記憶はどちらも記憶の節約だが、前者が作ったものの再利用、後者が呼び出す記憶の選択という違いがある。

深さ方向の計算を絞る発想もある。層ごとに上位選択で参加するトークンを絞り、系列と深さにわたって演算量を不均一に配る方式である\autocite{mod}。専門家を振り分ける仕組みとは振り分けの対象が分類的に異なる。早期に抜ける工夫もある。層の脱落と早期終了の損失で学習し、浅い層の草案を残りの層で確かめる自己推測につなげる\autocite{layerskip}。適応計算の先例としては、内部に分類器を添えて動的に抜ける古典的な工夫を、歴史の区切りとしてのみ位置づける\autocite{deebert}。深さの節約は注意の疎化や記憶の削減と直交しうる。どこを見るか、どれだけ残すか、どれだけ深く進むかを別々に絞ることで、正確な計算を保つ工夫と構造を変える設計の組み合わせが見通せる。

\subsection{共有化の系譜をもう一度たどる}
共有化の出発点は、すべての照会頭に対して鍵と値の頭を一つに束ねる発想である。照会側の多様性は残したまま、記憶側だけを共有することで、キャッシュの置き場所を節約できることが示された\autocite{mqa}。この節約が得られる代わりに、失われるのは鍵と値の投影の個別性である。各頭が別々の射影で文脈を切り取る自由度は狭まり、共有された一組の記憶に依存せざるを得なくなる。共有化の得失は、記憶量の削減と表現の個別性の後退という言葉で理解するのが正確であり、性能の数値ではなく設計上の取捨選択として捉えるべきである。

この単純な共有を、照会頭の群として段階化したのが束共有の考え方である。照会頭をいくつかの群に分け、群ごとに鍵と値を共有することで、完全共有と完全個別の中間を取る\autocite{gqa}。群を細かくすれば個別性に近づき、粗くすれば記憶節約に近づく。この中間化自体は機構としては素朴だが、運用上決定的だったのは既存のチェックポイントからの移行可能性である。多頭で学習された重みを起点に、全体を一から学び直すことなく束共有型へ移せることが示され、配備済みの資産を捨てずに省記憶化できる道が開かれた\autocite{gqa}。新規学習の高負担を避けられることは、研究上の優劣以上に現場での採用可否を左右する。

共有化の延長線上には、頭を束ねるのではなく鍵と値自体を小さな潜在表現に圧縮し、その潜在を照会頭間で共有する方向がある。既存の強力な基盤から継続学習のみで移行し、構造変更を最小限に抑えながら疎な選択機構を載せる実例が報告されている\autocite{v32report}。この構成側の工夫に対して、復号カーネル側の最適化が相補的に用意され、両者が合わさって効率化の物語を構成している\autocite{flashmla}。ここで注意すべきは、入出力の動かし方を洗練させて厳密な注意を保つ仕事とは役割が異なる点である。タイル化と逐次正規化により高帯域記憶への読み書きを抑えつつ近似なく計算する手法は、記憶の中身を変えるのではなく記憶の動かし方を変える\autocite{flashattn}。共有化が何を保持するかを問い、厳密な入出力最適化がどう動かすかを問う。この区別を保つことが、系譜をたどる際の見通しを良くする。

\subsection{疎化の二つの顔}
疎化には、学習の段階から疎さを組み込む顔と、学習後の推論時に不要と見なされる記憶を落とす顔がある。両者はともに注意対象を減らすが、適用できる時点と背負う制約がまったく異なる。混同すると、学習可能な選択と事後的な削除の区別が見えなくなる。

前者の代表は、圧縮と選択と局所走査を階層的に組み合わせ、学習過程から一貫して疎な注意を訓練する設計である。装置に整合したカーネルとともに端から端まで学習される点が、推論時のみ疎化する改修とは本質的に異なる\autocite{nsa}。既存基盤の上に疎な注意を載せ、長文脈の学習と推論の効率化を掲げて導入された事例もあり、索引機構による候補付けと細粒度のトークン選択を共有潜在の下で動かし、継続学習で移行する構成が報告されている\autocite{dsaexp}\autocite{dsadebut}\autocite{v32report}。局所注意の参照点としては、滑走窓と束共有を組み合わせた実用的な基盤が重要であり、疎な大域設計が何に対して局所性を主張しているのかを測る物差しとなる\autocite{mistral7b}。学習時に疎さを組み込む道は、索引と局所性を訓練に含める代わりに、学習基盤と実装の作り込みを要求する。

後者の代表は、生成時に蓄積されたキャッシュから残すものを選ぶ削除の技術である。注意の価値が少数のトークンに偏る観察に基づき、重要度の高いトークンと直近のトークンを残す方策が提案されている\autocite{h2o}。また、系列先頭のトークンが注意を引き付ける受け皿として働く現象が知られており、先頭の受け皿と滑走窓を残すことで窓型注意の性能を回復できることが示されている\autocite{sinks}。さらに、層ごとに注意の異質性があることを利用し、層間でキャッシュ割当を不均一に配分する考え方があり、層を一様に扱わない予算配分の視点を提供する\autocite{pyramidkv}。これらはいずれも凍結済み模型に適用できる軽さを持つが、削除である以上、削ってはならないものを削る危険と常に向き合う。

両者の使い分けは運用条件で決まる。学習や継続学習に介入できるなら、索引と局所性を組み込んだ疎さを訓練する道が選べる。すでに配備された重みを変えられず、長い文脈に対処しなければならないなら、事後的削除が現実解となる。相互作用として見逃せないのは、削除側が先頭受け皿の制約を受ける点である。いかなる削除方策も受け皿トークンを考慮しなければ後退するという制約は、削除の自由度に対する普遍的な留め金である\autocite{sinks}。また、層ごとに予算を考える発想は、削除だけでなく共有や選択の設計にも波及する。学習時選択と事後削除は競合ではなく、適用層の異なる補完関係にある。

\subsection{記憶の三つの置き場所}
記憶をどこに置くかという問いは、重みとキャッシュと外部再利用と条件付き参照を分けて考えると整理しやすい。常駐する重みは配備時に載り、要求間で動かない。増殖するキーと値のキャッシュは要求ごとに伸び、文脈が長いほど置き場所を食う。外部化された接頭辞や条件付きで引く表は、必要に応じて呼び出され、持ち運びと再利用の対象となる。何が動き、何が留まるかを見ることが運用理解の核心である。

常駐と増殖の対比は顕著である。重みは一度載れば反復利用される静的な資産であり、移動の主体ではない。これに対してキャッシュは生成の進行とともに追加され、保持と転送の負担が増す。厳密な入出力最適化が帯域の動かし方を改善しても、増殖自体を止めるわけではない\autocite{flashattn}。共有化や潜在圧縮が保持量を抑えるのは、この増殖の係数を小さくする働きである\autocite{mqa}\autocite{gqa}\autocite{flashmla}。層ごとの異質性を踏まえた予算配分が効くのも、増殖が層ごとに一様でないからである\autocite{pyramidkv}。

外部再利用は、同じバイト列を再び計算せずに使い回す発想である。ブロック単位の照合により要求をまたいでキャッシュ片を再利用する仕組みは、粒度をブロックに定めた再利用の具体化である\autocite{vllmprefix}。さらに、キャッシュの再利用や退避、分散を特定の推論基盤から切り離し、外部層として担う設計は、再利用と退避をまたぐ構成可能性を示す\autocite{lmcache}。再利用は計算の省略であり、接頭辞が一致すれば待たずに使える。選択がその都度の注意対象を選び取る判断であるのに対し、再利用はすでに作られたものを運んでくる手配である。このため、再利用の成否は一致率と粒度、配置と転送の設計に依存し、模型の賢さとは別の軸で効く。

条件付き参照は、記憶をすべて展開するのではなく、必要な断片だけを引く発想である。字句圧縮や複数頭のハッシュ、文脈に応じた調節、複数分岐の統合により、静的な表を参照計算として活用する設計が提案されている\autocite{engram}。公式実装が公開され、配置や退避の振る舞いに実行可能な裏付けが与えられている\autocite{engramimpl}。ホスト側に置かれた大規模な表を引く層構成の事例も報告されており、表をどこに置き何を移動させるかという配置設計の重要性を示す\autocite{qwenreadme}。ただし、この外部表と条件付き参照の設計との導出関係は確立しておらず、並行あるいは未解決として扱い、計算対記憶の圧力への収束的な応答と見るのが正確である\autocite{qwenlineage}。回帰的記憶の親に当たる門付き差分更新の系譜や、それを発展させた注意設計の報告も、この文脈に連なる\autocite{gdn}\autocite{kimilinear}\autocite{kda}。

結局、再利用と選択の区別は運用の言葉で言い直せる。再利用は当たれば得をする仕組みであり、外れても正しさは損なわれない。選択と削除は外すと品質に響く判断であり、受け皿や重要トークン、層別予算といった制約とともに設計しなければならない\autocite{vllmprefix}\autocite{lmcache}\autocite{h2o}\autocite{sinks}\autocite{pyramidkv}。常駐する重みを厚くするか、増殖するキャッシュを痩せさせるか、外部に置いて運ぶか、引くべき断片だけ引くか。この置き場所の設計こそが、注意と記憶を減らす努力の実質である。

\end{multicols}

\begin{center}
\small
\begin{tabularx}{\textwidth}{@{}p{0.22\textwidth}p{0.30\textwidth}X@{}}
\toprule
節約の位置 & 代表的な仕組み & 厳密か構造変更か \\
\midrule
鍵・値の共有化 & 単一頭共有・束共有 & 構造変更（保持を削る） \\
入出力の最適化 & タイル分割と逐次正規化 & 厳密（計算は変えない） \\
疎な選択 & 索引による注目先の限定 & 構造変更（見る場所を変える） \\
局所窓 & 窓内のみの注意 & 構造変更（大域を捨てる） \\
線形・回帰記憶 & 門付き差分と交互配置 & 構造変更（二次計算を置換） \\
保持の除去 & 重要度・吸着・層配分で残す & 構造変更（残す量を絞る） \\
再利用と条件付き記憶 & 接頭辞共有・外部層・選択参照 & 厳密な再利用と記憶の選択 \\
深さの適応 & 層別参加・早期終了・自己推測 & 構造変更（進む深さを絞る） \\
\bottomrule
\end{tabularx}
\end{center}
\autocite{mqa,gqa,flashattn,nsa,kimilinear,h2o,pyramidkv,engram}

\begin{claimboundary}[本節の読解上の境界]
由来の未確定な系譜は断定しない。呼称の衝突が未解決の別系列には触れない。数値表・核実装・速度の事実は引用せず、節範囲と告知範囲の読みに閉じる。
\end{claimboundary}

\section{Bitを減らし、巨大モデルを手元で動かす}
\label{sec:precision}
\sectionkicker{PRECISION AND LOCAL}

\begin{multicols}{2}
\raggedcolumns
巨大な言語モデルを手元で動かすとき、最初の壁は重みの置き場所である。初期の有力な答えは、ごく少数の特徴次元に現れる外れ値だけを高精度計算に残し、残りの大部分を八ビット整数で処理する分解だった\autocite{llmint8}。外れ値次元は全体のごく一部でありながら量子化誤差を支配するため、そこだけを丁寧に扱うことで記憶容量を大きく減らしつつ精度劣化を抑えられる。続く一回性の重み限定の学習後量子化は、近似的な二次情報を手がかりに層ごとの再構成誤差を抑え、三から四ビット級の重みにまで圧縮しながら困惑度の増加を小さく保つ道を示した\autocite{gptq}。追加学習を前提とせず、少量の較正データで済む点が実務上の意味を持つ。

ここで分類を明確にしたい。重みだけを削る手法と重みと活性値の双方を削る手法は、誤差の出所も対策も異なる。さらに、学習後に合わせる手法と学習時から低ビットを前提にする手法、微調整時に量子化する手法は別の系統であり、同一の物差しで優劣を論じられない。鍵値キャッシュの量子化や最適化器の精度、保存符号化も重み量子化と同一視できない別軸である。詳細な外れ値の出現分析や層別誤差の導出には立ち入らず、概要水準の主張として扱う。

\subsection{同時量子化の二つの道}
重みと活性値の双方を八ビット化する道は、重みだけを削る道とは別の難しさを持つ。平滑化の発想は、活性値の振れを経路ごとの係数で均し、対応する分だけ重みを調整することで、再学習なしに重みと活性値の八ビット化を可能にする\autocite{smoothquant}。係数決定の数式導出には立ち入らず、難所の移行という要点に絞って扱う。一方、活性値統計から重要な重みを守る発想は、ごく少量の較正集合でも強い四ビット化を実現する\autocite{awq}。すべての重みを一様に削るのではなく、出力への影響が大きい部分を保護し、残りに厳しい量子化を割り当てる。この重要度保護の考え方は、手元の実行環境で用いられる混合量子化や重要度行列の実践へと連なる。平滑係数の導出過程や顕著性分析の細部は摂取範囲外であり、機構の要点を概要水準で扱う。

\subsection{器と削り方を分ける}
手元で巨大モデルを扱う議論では、保存形式と量子化手法の混同がしばしば起きる。\textbf{GGUFは量子化手法ではなく表現／コンテナ形式である}\autocite{gguf}。この形式は、テンソル配置や付随情報の書式、記憶写像で読み出せる搭載方法など、局所推論の生態系が共有する器を定める。一方、重要度に応じた混合量子化や動的量子化の系列は、その器に載せるテンソルへ適用される手法であり、形式自体の性質ではない。同じ器に入っていても、中身の削り方が違えば振る舞いは変わる。この区別が、容量と品質の議論の前提になる。

参照実装としての局所推論実行系は、この形式の読み書きと多様な量子化系列、中央演算装置と画像演算装置への分割実行を一体で担い、新規モデルの初期対応の受け皿になってきた\autocite{llamacpp}。層の一部を装置間で分担し、収まらない分を主記憶側に置く運用が、手元再現の現実的な姿である。公開文書で示される大容量モデルの器の用意は、手元の記憶容量に収まる選択肢の案内としてのみ扱う\autocite{unslothglm}。品質指標や多段予測による高速化の数値は提供元による測定であり、独立検証が未了のため事実として採用せず一般化しない。

\subsection{浮動小数点の低ビット化}
浮動小数点の低ビット化は、整数化とは別の設計軸である。八ビット浮動小数点の二形式の組み合わせは、指数部と仮数部の配分が異なる二者を用途に応じて使い分ける混合精度の方法論として提示された\autocite{fp8formats}。さらに、微小スケーリングの公開仕様は、小さな要素群ごとに共有指数を備えることで、狭いビット幅でも精度を保つ枠組みを与える\autocite{mxspec}。群幅三十二での共有スケールと、四ビット要素に共有指数を組み合わせた取り決めが、その具体例である。この系列の配備例として、高密度部を八ビット微小浮動小数点に置き、専門家部を四ビット微小浮動小数点に置く検査点構成が、文書化された手順の水準で報告されている\autocite{mxfp4}。鍵値キャッシュの量子化や最適化器精度とは軸が異なることに注意したい。配送符号の査読は行っておらず、実装正当性の断定はしない。文書化の水準を超えて語らない。

\subsection{適応と原生低ビット}
微調整の効率化は、量子化と適応の接点にある。凍結した四ビット基盤に低ランク付加器を組み合わせる方法は、大規模モデルの振る舞いを保ったまま学習対象を小さくし、単一の画像演算装置でも扱えるようにする\autocite{qlora}。ここで最適化器の精度や記憶管理は、重み自体の量子化とは別の軸であり、混同できない。一方、学習時から低ビットを前提にする原生低ビットの方向は、学習後に削る発想とは系統が異なる。一ビットで学習する変圧器の試みは、十分な規模で高精度に匹敵する振る舞いを示しうると報告される\autocite{bitnet}。三値重みによる一・五八ビットの試みは、数十億規模で同等の困惑度や課題性能を保ちつつ遅延や記憶、流通量、電力の利点を報告する\autocite{bitnet158}。ただし、特定企業を離れた産業配備は未確認であり、方向性の証拠として扱い配備の証拠としては扱わない。

\subsection{総量と作動量の落差}
巨大な混合専門家モデルを手元で動かす現実は、総量と作動量の落差にある。不均質実行の枠組みは、この性質を利用し、中央演算装置と画像演算装置に専門家を配置換えしながら必要分だけ読み出すことで、容量と計算を切り離す\autocite{ktransformers}。総量が大きくても作動量が小さければ手元の可能性が開ける。最新の旗艦級混合専門家モデルの手元移植は、変換器や疎注意、予測表などの構成を含む実装報告として公開され、検証が報告者側の記録として示されている\autocite{qwenpr}。予測表の主記憶挙動も実装記述の水準であり、独立した再構築や測定は行っておらず、正当性の断定はしない。実践者の受容記録は、配置に依存する観察として扱い、検証済み事実とはしない。技術的事実や経済的事実、許諾条件の事実をこの記録に依拠して述べることはしない。

\subsection{外れ値と重要度の二つの見方}
巨大モデルの精度を保ったままビット数を減らす試みが、一様な丸めではうまくいかない理由は、誤差が均等に広がるのではなく、ごく一部の場所に集中するからである。一部の特徴次元には突出して大きな値が現れ、一部の重みは活性化の振る舞いと結びついて強く保護に値する。どこを丁寧に扱うべきかを見極めれば、全体を高精度に保つ必要はなくなる。この見極めに二つの流儀がある。一つは外れ値そのものを分離する考え方であり、もう一つは活性化の統計から重要度を読み取る考え方である\autocite{llmint8}\autocite{smoothquant}\autocite{awq}。

外れ値駆動の分解という流儀を切り開いたのが、突出した外れ値次元だけを高い精度に残し、残りの大部分を低いビット幅で計算する手法である\autocite{llmint8}。全体を一律に低ビット化すれば、外れ値が量子化範囲を引っ張り、通常の値の分解能が犠牲になる。そこで計算の経路を分け、少数の扱いにくい次元だけを別扱いにする。こうした分離によって、制約の強い記憶容量でも巨大なモデルの推論が可能になるという方向性が示された。ここでの発見は単なる実装の工夫にとどまらず、後に続く手法が共有する前提となった。すなわち、活性化に現れる外れ値こそが量子化の難しさの中心であり、その扱いをどう設計するかが手法の性格を決めるという見方である。

もう一つの流儀は、活性化の統計を手がかりに、守るべき重みを見つける考え方である。重みだけを低ビット化する学習後量子化の枠組みでは、層ごとの誤差を二次情報を用いて補償する手法が代表的であり、学習し直すことなく低いビット幅へ到達できることが示された\autocite{gptq}。これに対して、重みと活性化の双方を低ビット化しようとすると、活性化側の外れ値が新たな壁になる。この壁を、チャネルごとの平滑化によって重み側へ難しさを移し替えることで乗り越えようとしたのが、重みと活性化の同時量子化を学習なしで可能にする手法である\autocite{smoothquant}。平滑化は、外れ値の激しいチャネルを穏やかに均し、その分だけ対応する重み側に負担を寄せる操作として理解できる。活性化の量子化を容易にする代償を、比較的対処しやすい重み側で引き受けるという分担の発想である。

活性化から重要度を読み取る方向をさらに推し進めたのが、顕著な重みを保護する手法である\autocite{awq}。この手法は、活性化の統計を観察することで、どの重みがモデルの振る舞いを強く左右するかを判定し、その重要な重みが量子化誤差で損なわれないよう優先的に守る。重要な少数を厚く守り、それ以外を思い切って削るという配分が、低いビット幅でも精度を保つ鍵になる。この重要度保護の考え方は、後に手元での実行環境における混合量子化や重要度行列の実践へと直接つながっていく。較正データの倹約性もここから生まれる。必要とされるのは大規模な学習用データではなく、モデルの振る舞いの癖を読み取るための少量の文章である。学習し直しを要せず、手元の少量の観察だけで適用できることは、手元の装置との相性を決定的に良くする。こうした軽さが、後の共有容器と結びついて、手元での試行錯誤を広げた。

二つの見方がともに一様な削減を上回るのは、精度の予算を誤差の痛い場所へ集中投下するからである。外れ値の分離は痛みの発生源を隔離し、平滑化は痛みの置き場所を移し替え、重要度保護は痛みに弱い箇所を事前に補強する。いずれも、全体を同じ粗さで扱うのではなく、観察された振る舞いに応じて扱いを変えるという点で共通している。

\subsection{形式の標準化がもたらしたもの}
手元で巨大モデルを動かす実践を支えたのは、個別の量子化手法の巧妙さだけではなく、それらを運ぶ器の共通化である。コンテナ形式は、テンソルの配置や付随するメタ情報、読み出しのための載置の仕方を定め、実行環境が同じ手順で多様なモデルを開けるようにする\autocite{gguf}。メモリへの写像によって大きな重みを効率よく扱い、中央演算装置と画像処理装置の間で処理を分担する実行形態と組み合わせることで、容量の大きな重みと実際に計算する部分とを切り離す運用が可能になった\autocite{llamacpp}。この器の上に、量子化の系列が載せられる。特定の量子化の系列は器に格納される対象であり、器そのものの性質ではない。ここで改めて確認しておきたいのは、GGUFは量子化手法ではなく表現／コンテナ形式であるということだ。

共有された器がもたらした第一の変化は、比較可能性である。同じ容器、同じ読み手、同じ実行手順の上に異なる量子化が載るとき、利用者は手法の違いを同じ土俵で比べられる。変換器や読み込み手順、分割実行の仕組みを共有できることは、手順の共有という第二の変化を生む。新しいモデルが現れたとき、器と手順が共通であれば、移植の作業は一から作り直す仕事ではなく、差分を埋める仕事になる。実際に、混合エキスパート型の大規模モデルの手元への移植は、変換器やグラフ、検証手順の報告を伴いながら、参照実装の上に積み上げられる形で進められた\autocite{llamacpp}\autocite{qwenpr}。派生した実行系による文書化や配布の整備も、同じ器を前提とするからこそ利用者に届きやすくなる\autocite{unslothglm}。

第三の変化は、異種配置という運用の定着である。総量の大きな重みを手元の記憶容量に収め、実際に働くごく一部だけを高速な装置で計算するという分離は、混合エキスパート型の重みと記憶の問題に対する実践的な答えとなった\autocite{ktransformers}。重みの置き場所と計算の場所を切り離す考え方は、容器による写像や分割実行の仕組みと相性が良く、容量と速度の異なる装置を組み合わせる運用を後押しした。仕様が器を定め、手法が載せ物を定めるという分担が、比較と共有の循環を回したのである。

\subsection{低ビットの三つの時間軸}
低ビット化をいつ行うかという問いは、誤差がどこで生まれるかという問いと表裏一体である。学習が終わった後に丸めるのか、適応のための学習と組み合わせるのか、それとも学習そのものを低いビット幅で設計するのかによって、誤差の源も、対処の姿も変わる。この時間軸の違いを見分けることが、手元で動かすための選択を誤らないための助けになる\autocite{gptq}\autocite{qlora}\autocite{bitnet}\autocite{bitnet158}。

第一の時間軸は、学習後の量子化である。学習済みの重みを、学習し直すことなく低いビット幅へ変換する立場であり、重みのみを対象とする一回的な変換がその代表である\autocite{gptq}。ここでの誤差は、事後的な丸めから生まれる。対処は、層ごとの誤差を補償する工夫、外れ値の扱いを変える平滑化、重要な重みの保護といった、観察と補償の技術になる\autocite{smoothquant}\autocite{awq}。利点は適用の軽さであり、手元の少量の観察で済み、学習基盤を持たない利用者でも試せることである。限界もまた明確であり、学習過程そのものに手を入れない以上、丸めによって失われた情報を取り戻すことはできない。

第二の時間軸は、量子化と適応を組み合わせる立場である。低いビット幅で凍結した土台の上に、小さな学習可能な付加部分を載せて調整する手法がその代表であり、量子化の選択が微調整の効率を直接左右することを示した\autocite{qlora}。ここでの誤差は、単なる丸めの誤差ではなく、限られた表現力の上でどれだけ目的の振る舞いに寄せられるかという適応の誤差を含む。土台を軽くすれば手元の装置でも調整が可能になり、適応の自由度を残せば用途への合わせ込みができる。この時間軸は、量子化と適応という別々の営みが、精度の予算を通じて結びついていることを教える。削り方の選択が、学び方の効率を決めるのである。

第三の時間軸は、はじめから低いビット幅で学習する立場である。二値や三値の重みで大規模な変換器を学習し、その規模拡大の振る舞いを全精度と対比可能なものとして示そうとする方向がこれにあたる\autocite{bitnet}\autocite{bitnet158}。ここでの誤差は事後的な丸めではなく、低い表現力のなかで学習が安定に進むかという学習そのものの難しさから生まれる。対処もまた、学習手順や規模則、表現形式の設計になる。微小な粒度で共有縮尺を定める形式の標準化や、少数の要素形式の整備は、この方向を支える土台として理解できる\autocite{fp8formats}\autocite{mxspec}。実際に、専門家ごとに異なる精度を割り当て、密な部分や埋め込みや出力層の扱いを分けた構成が、運用の現場で用いられる事例として報告されてきた\autocite{mxfp4}。ただし、研究としての方向性の証拠と、広く配備された実績とは区別されねばならない。後者については開かれた問いとして残すのが、証拠の範囲に忠実な態度である。

さらに見落としてはならないのは、重みのビット幅とは別の軸が二つあることである。一つは学習や適応のための記憶であり、最適化の状態や中間結果が占める容量は、重みそのものの軽さとは別に考えねばならない。もう一つは、生成の過程で蓄積される鍵と値の記憶であり、その量子化は重みの量子化とは異なる誤差の振る舞いを持つ。重みを削ること、学習の記憶を抑えること、生成時の記憶を抑えることは、それぞれ誤差の源も対処も異なる別問題である。三つの時間軸と二つの別軸を混同せず、どの軸のどの誤差に対処しているのかを問い続けることが、ビットを減らして手元で動かす試みを読み解くための確かな物差しになる。

\end{multicols}

\begin{center}
\small
\begin{tabularx}{\textwidth}{@{}p{0.20\textwidth}p{0.30\textwidth}X@{}}
\toprule
削る対象と時点 & 代表的な手法 & 分離すべき軸 \\
\midrule
重み限定・学習後 & 外れ値分離・再構成誤差の抑制 & 活性値同時量子化と混同しない \\
重み活性同時・学習後 & 平滑化・重要度保護 & 重み限定と誤差の出所が異なる \\
微調整時の量子化 & 凍結基盤と低ランク付加 & 最適化器精度は別軸 \\
原生低ビット & 一ビット学習・三値重み & 学習後量子化の代替と断定しない \\
浮動小数点の低ビット & 混合精度・微小スケーリング & 整数化と別の設計軸 \\
保存の器 & 表現・コンテナ形式 & 量子化手法ではない \\
局所実行 & 分割実行・不均質配置 & 総量と作動量を分ける \\
\bottomrule
\end{tabularx}
\end{center}
\autocite{llmint8,gptq,qlora,gguf,fp8formats,bitnet}

\begin{communitynote}[実践者の受容---配置依存の観察]
手元配備の可否や外部退避と主記憶と帯域が隘路になること、量子化選択と実行系の成熟が使い勝手を左右することは、実践者の報告として共有されている\autocite{qwenpr,unslothglm}。ただし配置に依存する観察であり、検証済み事実とはしない。技術的事実や経済的事実の根拠には用いない。
\end{communitynote}

\begin{claimboundary}[本節の読解上の境界]
GGUFは表現・コンテナ形式であり量子化手法ではない。提供元測定の品質・速度数値は隔離し、独立検証まで一般化しない。符号庫の査読や再構築は行わず、産業配備の証拠としては扱わない。
\end{claimboundary}

\section{1 tokenずつ待たない}
\label{sec:decode}
\sectionkicker{DECODING AND BUDGET}

\begin{multicols}{2}
\raggedcolumns
生成モデルの待ち時間の多くは、一トークンずつ順番に確定させなければならない直列性に由来する。投機的デコーディングはこの構造を正面から置き換えるのではなく、小さなドラフト側が複数トークンを先に提案し、標的モデルがそれらをまとめて検証する手順に変える\autocite{specdecode,chenspec}。受け入れられた接頭辞だけを確定し、外れた位置からは標的分布に従って引き直すため、出力の分布を保ったまま無駄な逐次呼び出しを減らせる。標的モデルの再学習を要しない点も重要で、既存の重みをそのまま検証器として使える。速度の鍵は受容率であり、ドラフト案がどれだけ標的に受け入れられるかが全体の効率を決めるという整理が、以後の研究を貫く評価軸になった。ただし受理条件の厳密な証明や速度表の全体は本稿の消費範囲に含まれず、機構と保証の骨格に限って扱う。\textbf{直列性の緩和と推論予算の節約は隣り合うが別問題である}。

\subsection{先読みの質を上げる}
ドラフトの質を上げる方向で転機になったのが、特徴量レベルで先読みする方式である\autocite{eagle}。トークン表層ではなく中間特徴量から次の展開を予測し、単層の軽量復号器と木構造の注意機構で複数候補を束ねて検証する。基幹モデルの微調整を要せず出力分布を変えない無損失の範囲が保たれ、量子化やコンパイルと組み合わせ可能とされる点も実務上の含意が大きい。先行する多頭方式は、凍結した基幹モデルに軽量な複数ヘッドと木検証を付加し、別個のドラフトモデルを不要にした前例として位置づけられる\autocite{medusa}。貪欲復号での無損失性は保たれる一方、非貪欲生成の品質は限定付きの対象であり、無条件の保証ではない。発展形では木を文脈依存に動的化し、受理モデルの予測で展開を変える\autocite{eagle2}。ただし木構造や学習データのアブレーションは部分的であり、単一の測定条件の数値を普遍定数として扱うことはできない。

複数トークンを同時に学ぶ訓練が、推論の先読みに転用される流れも押さえておきたい。共通胴体の上に独立な複数出力ヘッドを置き、次トークン予測に加えて先のトークンを補助課題として学ぶ方式では、訓練時の追加負担なしに標本効率が上がり、規模が大きいほど効果が伸びると報告されている\autocite{mtp}。現在の実装面では、推論サーバの文書において標的モデル自身の多トークン層が草案を作る方式が案内され、別個のドラフトモデルを要しない運用が示されている\autocite{vllmmtp}。ただし文書本文の全体は未取得であり、主張は要約と下流の裏づけに限定される。ベンダ側の速度数値は隔離し、独立の確認が済むまでは引用しない。

\subsection{足し算の効率}
テスト時計算の配分に目を向けると、大きなモデルに置き換える前に考えるべき選択肢が見えてくる。難易度適応的な検証器探索と分布改訂を組み合わせた最適配分では、素朴な最良候補選択の反復に対して効率改善が本文の要旨として報告されている\autocite{snell2024}。検証器の訓練の詳細は未取得であり、万能の優劣として読むのではなく、基礎能力が残る領域での条件付きの効率論として受け止める必要がある。単純なテスト時拡張の事例では、数万件から千件への難易度と多様性と品質による精選が土台になる\autocite{s1}。思考を打ち切ったり末尾に待機語を付して延長したりする予算強制が核になり、短すぎる思考を延ばして推論を修復する操作として働く。予算の付け方が結果を左右するという点で、計算資源の配分は模型選択と同等の設計変数である。

\subsection{引き算の効率}
長い思考を短く保つ学習的な試みでは、制限を報酬設計に織り込む発想が採られている\autocite{thinkprune}。制限を越えた出力は報酬計算の前に破棄される厳格な上限のもとで強化学習を行い、上限を段階的に締めながら要点を選び直す。本文の結果表では、蒸留由来の小規模学習対象において正答率をほぼ保ったまま生成長がおおむね半減する様子が示されている。温度や試行回数や上限長を明示した手順とともに示され、評価条件の束縛が明示されている点が重要で、汎用的な短縮保証ではない。発見的な打ち切りではなく学習された方策として位置づけられ、重みを更新する手法と推論時の置換策の境界が保たれている。背景には大規模強化学習で推論能力を引き出し小型高密度模型へ蒸留する流れがあり\autocite{r1}、対象外の規模や数学外課題への一般化は未確定である。訓練費用の会計は抽出されておらず、推論効率の改善と学習費用の総量は切り分けて読む必要がある。

考えさせれば考えさせるほど良いわけではないという観測が、予算設計の現実味を支えている\autocite{overthink}。長い推論では正答率が上がってから下がる山形の推移が見られ、正しかった答えを捨ててしまう反転事象が記録されている。計算量を罰則付きの効用で捉える枠組みでは、費用に敏感な条件では早期停止が同等の精度で大きく節約するという整理が示されている。例示の算術を請求書の実額と混同せず、費用対効果の考え方として読むことが重要である。運用界面では、数値で推論努力を指定する系列と、思考の可否や段階離散の強弱で指定する系列という、明示的予算の操作系が報告されている\autocite{effortbind}。ただし界面文書の全体は未取得で、段階意味の把握は要約水準に留まる。復号の待ちと思考の浪費を混同せず、層の異なる無駄として設計する視点が求められる。

\subsection{受容率という単一の物差し}
投機的デコーディングの成否を決めるのは、草案の巧拙でも検証の厳しさでもなく、両者の噛み合わせが受容率という一つの量に畳み込まれた結果である。草案が先に複数候補を並べ、標的模型がそれをまとめて検証し、受け入れられた接頭辞だけを残す。この往復では、草案がどれだけ速く巧みに書けても、標的の分布がそれを撥ねつければ前進は生まれない。逆に、標的の判定が甘ければ速く進むが、守るべき分布そのものが崩れる。だから受容率は単なる中間指標ではなく、速度と正しさの交換を一本化する物差しとして残り続ける\autocite{specdecode}。

この枠組みに厳密な形を与えたのが、受容判定の定式化である。草案分布と標的分布のずれを補正する受容条件を明示することで、分布を保ったままの加速が何によって保証されるのかが見えるようになった\autocite{chenspec}。無制約な近似ではなく、受容条件を通った分だけが標的分布の再現として正当化される。この条件があるからこそ、後の手法は貪欲生成の厳密さを保つ部分と、非貪欲生成で品質保証の範囲を限定せざるを得ない部分とを切り分けて議論できる。

草案の質と標的の厳しさは相補的であるという見方はここから自然に導かれる。軽量な複数頭で基幹を凍結したまま草案を出す方式は、別立ての草案模型を持たずに検証の木構造で受容の機会を増やす発想である\autocite{medusa}。貪欲生成の範囲では分布の不変性が保たれる一方、非貪欲生成では品質の主張は限定付きになる。この限定は弱さではなく、受容率の物差しで見た正直な適用範囲の宣言である。特徴量水準で草案を作る方式は、この相補性をさらに押し進め、表層の語彙予測ではなく内部表現の連続性に草案の根拠を求める。本体を微調整せずに出力分布を変えない立場は、量子化やコンパイルといった実装上の最適化と組み合わせやすいことも含めて、受容率を下げないための設計として読める\autocite{eagle}。

草案の木を文脈に応じて動的に組み替える考え方は、受容率を事後的な平均値ではなく、文脈ごとの予測対象として扱う転換である\autocite{eagle2}。ただし、この物差しの適用範囲には開いた問いが残る。とりわけ、バッチが大きくなったときや装置構成が変わったときに受容率と実効速度の関係がどうずれるかは、現時点の証拠では確定していない。受容率が高いことと運用条件で速いことが常に一致するとは限らない。受容率を不変の真理ではなく、条件付きの物差しとして扱う姿勢を保つべきである。

\subsection{学習の改善としての先読み}
複数トークン予測は、しばしば復号の工夫としてだけ語られるが、その出自は学習の改善である。共有された基幹部に独立な複数の出力頭を付け、未来の複数位置を同時に予測させる補助課題は、学習時の負担を実質的に増やさずに標本効率を高め、規模が大きくなるほど効果が育つ方向に働く\autocite{mtp}。速く生成するための仕掛けを後付けするのではなく、未来を見越して学ぶことがそもそも良い模型を作るという順序の逆転にある。

実装面での帰結は、標的模型自身の先読み層が草案を担い、別立ての草案模型を不要にする形態である。文書化された実装経路は、要約水準の証拠に裏打ちされた範囲ではあるが、主要な系統が先読みを目的関数に取り込む流れと整合的である\autocite{vllmmtp}。旗艦的な系統における採用を、単なる高速化の流行としてではなく、目的関数の出自として読むことが、この小節の狙いである。学習時に未来を予測させた模型は、復号時に未来を草案できる。推論基盤がその能力を直接呼び出せる形で整備されたとき、先読みは研究の工夫から運用の前提へと位相を変える。

この整理は、推論様式の系譜にも接続する。強化学習で育てた推論様式を小さな生徒模型に移す流れは、現在の利用の錨として残る\autocite{r1}。先読みの採用史も同様に、個別の高速化手法の列挙としてではなく、どの世代が何を目的関数に据えたかという系譜として読むべきである。学習を良くすることが復号の選択肢を増やし、復号の選択肢が学習の価値を運用に変える。この循環こそ、先読みを単なる先回り生成と呼ばず、学習の改善として扱う理由である。

\subsection{予算表で考える}
テスト時計算を一枚の予算表として眺めると、計算を足す方向と削る方向は対立する主張ではなく、同じ表の左右に見える。足す方向の代表は、検証器付き探索や分布の改訂を難易度に応じて配分する考え方である。難しい問いには探索を厚くし、易しい問いには生成の手直しで済ませる。この適応的な配分が、単純な多数決的な基準を上回る効率を生むというのが、計算最適という言葉の実質である\autocite{snell2024}。検証器の学習の細部は未確定だが、難易度適応という濃淡の存在自体は揺らがない。

同じ表の足す方向のもう一つの姿が、思考の打ち切りと強制延長である。厳選された少量の高難度・多様・高品質な問いに対して、思考を途中で切ることも、待機を促す合図で延ばすことも、どちらも予算操作として同列に扱える。量や多様性や難易度をめぐる詳細な切り分けは構造的水準にとどまるが、予算を外から与えるという操作の素朴さと、その効果の堅牢さが要点である\autocite{s1}。

削る方向の代表は、厳しい上限のもとで強化学習により短く考えることを学ばせる方法である\autocite{thinkprune}。上限を超えた出力は正誤にかかわらず報酬計算から除外されるため、模型は限られた長さに収まる解法を獲得せざるを得ない。この手法が依拠する推論蒸留の系譜は、強化学習で育成した推論様式を生徒系統に移す流れに属する\autocite{r1}。

削る方向を正当化するもう一つの視点が、考えすぎの分析である。長い推論は常に良いわけではなく、正しかった答えを途中で捨てる反転を伴いながら、伸ばすほど精度が頭打ちから低下に転じることがある\autocite{overthink}。反転機構の因果分析は部分的であり、費用の算術は請求書ではなく例示だが、予算表に費用の欄を設けるべきだという主張は明確である。運用界面における明示的な予算つまみ、すなわち数値や段階で推論の努力量を結び付ける機構は、この費用意識の生産形態である\autocite{effortbind}。

予算表を閉じるときに欠かせないのが、学習費用と推論費用の分離である。強化学習による短縮の学習費用は抽出されておらず、推論時の節約と混同してはならない。読者が持つべき表は、行に難易度、列に足す操作と削る操作、欄外に学習費用と推論費用を置いたものである。復号の直列性と推論予算の浪費は同じ機構ではなく、隣り合う課題であることを忘れずに、この表を運用の道具として使ってほしい。

\end{multicols}

\begin{center}
\small
\begin{tabularx}{\textwidth}{@{}p{0.18\textwidth}p{0.26\textwidth}p{0.26\textwidth}X@{}}
\toprule
予算表の軸 & 足す操作 & 削る操作 & 欄外の費用 \\
\midrule
難易度の配分 & 検証器付き探索と分布改訂 & 打ち切りと強制延長 & 検証器の育成費用は別勘定 \\
草案の質 & 受容率の向上と動的木 & 先読み学習の転用 & 別立て模型の不要化は運用の利得 \\
思考の長さ & 待機合図による延長と立て直し & 上限制約の学習と早期停止 & 短縮の学習費用は未抽出 \\
努力量の指定 & 数値と段階のつまみ & 中程度での回収 & 界面仕様は要約水準 \\
分布の保存 & 受容条件の統制 & 置換策との混同の回避 & 重み更新と置換は別機構 \\
\bottomrule
\end{tabularx}
\end{center}
\autocite{specdecode,eagle,mtp,snell2024,s1,thinkprune,overthink,effortbind}

\begin{center}
\small
\begin{tabularx}{\textwidth}{@{}p{0.22\textwidth}p{0.32\textwidth}X@{}}
\toprule
待ちと予算の設計 & 仕組みの要点 & 隣接するが別機構 \\
\midrule
投機的検証 & 小さな草案を標的が一括検証 & 分布保存と受容率が鍵 \\
特徴量の先読み & 中間特徴量から複数候補を束ねる & 表層予測の限界を越える \\
先読みの学習 & 複数トークンを補助課題として学ぶ & 学習の改善と復号の再利用 \\
検証器つき探索 & 難易度適応と分布改訂 & 基礎能力が残る条件付き \\
予算強制 & 打ち切りと延長で思考を修復 & 短すぎる思考の延ばし \\
学習的短縮 & 上限制約の強化学習で要点化 & 重み更新と置換策の境界 \\
過思考の知見 & 山形推移と反転の記録 & 止める判断が精度維持 \\
明示的予算 & 努力量の数値指定と段階指定 & 要約水準の把握 \\
\bottomrule
\end{tabularx}
\end{center}
\autocite{specdecode,eagle,mtp,snell2024,thinkprune,overthink}

\begin{claimboundary}[本節の読解上の境界]
証明の細部、速度表の全体、アブレーション、訓練費用の会計、対象外規模への一般化には立ち入らない。ベンダ側の速度数値は隔離する。直列性と予算浪費を一括りに語らない。
\end{claimboundary}

\section{Servingで消える無駄}
\label{sec:serving}
\sectionkicker{SERVING EFFICIENCY}

\begin{multicols}{2}
\raggedcolumns
大規模言語モデルの運用で最初に詰まるのは、重みの大きさではなく生成途中に膨らむ記憶である。ここでの転換点は、鍵と値の履歴を連続領域として抱え込まず、小さな塊に分けて必要な分だけ結び付けて管理する発想である\autocite{pagedattn}。塊単位になれば再利用や受け渡しの単位も揃い、空きの再利用が細かく回る。到着順に束ね直す連続的な束ね方は、この記憶管理があって初めて安定する。その後の接頭辞再利用や分離構成、外部への退避はいずれも、履歴を塊として照合し移動できることを前提にしており、この層の有無が後段の工夫の効きを決める。本稿では詳細な表設計や条件付き測定の全体には踏み込まず、確立した機構としての役割に限って扱う。伝えたいのは、\textbf{構造が巧妙なモデルでも記憶の管理と束ね方が粗ければ運用の効率は伸びない}という順序である。

\subsection{重なりの使い回し}
対話や道具呼び出しを重ねる使い方では、同じ前置きを何度も計算し直す無駄が膨らむ。ここでの工夫は、接頭辞の共有関係を木のように捉え、一致する部分の計算結果を使い回すことである\autocite{sglang}。制約付き生成のための状態管理を圧縮して持つ発想も、同じ枝刈りの延長にある。さらに、再利用の仕組みを特定の実行機構の内部に閉じ込めず、外部の層として切り出して複数の機構から使えるようにする発想がある\autocite{lmcache}。使い回しの設計は応答の速さだけでなく費用の読みやすさにも効く。粒度の違いも重要で、塊の照合による再利用と木による管理は互いに含み合う関係ではなく、得意な共有の形が異なる。どちらか一方が他方を包摂するものとして語らず、使い分けとして理解すべきである。本稿では設計の細部や条件付き評価の全体には踏み込まず、繰り返し呼び出しの無駄を減らす確立した役割に限って述べる。

\subsection{段階の分離}
生成は前半と後半で性質が異なる。前半は文脈を一気に読み込む計算中心の段階であり、後半は一つずつ言葉を紡ぐ記憶帯域中心の段階である。両者を同じ資源プールで同じさばき方に縛ると、どちらかに合わせた調整がもう一方の無駄になる。ここでの発想は、段階ごとに適した資源の集まりを用意し、役割分担で全体をならすことである\autocite{splitwise}。ただし分離には代償がある。生成途中の記憶を受け渡す転送が新たに生まれ、特化の利益と受け渡しの費用の釣り合いが問われる。記憶を分散した共有資源として束ね、必要に応じて高速な転送で結ぶ考え方は、この釣り合いを運用する土台になる\autocite{mooncake}。本稿では分離の配置解析や費用表、分散保持の手順や評価の全体には踏み込まず、確立した役割の範囲で扱い、補完的な記録に基づく部分は抽象の範囲に留める。

同時処理では、長い入力の読み込みが短い応答の生成を邪魔するという噛み合わせの悪さが出る。読み込みを小さな断片に分けて生成の合間に滑り込ませるやり方が、その処方箋である\autocite{sarathi}。断片化すれば生成の遅れを抑えながら読み込みを進められ、全体の詰まりがほぐれる。後続の展開は、その釣り合いを運用条件に合わせて整える方向に進んだと位置づけられる。運用条件が変われば最適な区切り方も変わるため、一つの正解を固定せず調整の幅を持たせることが重要になる。本稿では分割手順の細部や得失の分析の全体には踏み込まず、補完的な記録に基づく確立した役割の範囲で述べる。

\subsection{核と通信と基盤}
運用の軽さは、待ち行列の工夫だけでは決まらない。圧縮された注意の復号に寄せた核は、構造側の工夫と対になる実行側の工夫であり、復号の重さを直接引き受ける\autocite{flashmla}。多数の専門家に分散する構成では、分散と結合の通信が全体の重さを左右し、通信のさばき方が計算の工夫と同じ重みを持つ\autocite{deepep}。行列積を細かな粒度で低精度化する土台は、大規模な学習の数値的な成立を支え、学習と運用をつなぐ\autocite{deepgemm}。複数の技法を一つの実行基盤に同居させる共有核の考え方は、個別の最適化を組み合わせて使えるようにする接着剤である\autocite{flashinfer}。核と通信と数値基盤の三層が揃って初めて、待ち行列の工夫が実を結ぶ。本稿では核の実装や数値解析、通信量の測定の全体には踏み込まず、確立した役割に限って述べ、個別の高速化数値は引用しない。

運用の仕上げでは、既存の先読み部品を外して別の推測の仕組みに置き換えるといった選択が現れる\autocite{dspark}。公開された導入手順は生きた文書であり、特定の固定構成での確認はその構成での測定に留まる\autocite{vllmrecipe}。他の画像や上限設定、異なる資源量への転用は保証されず、文書の更新に伴う揺らぎを前提に読むべきである。構造の効率と運用の効率は別の勘定であり、両者を混同すると投資の順序を誤る。まず漏れを塞ぎ、次に構造の利益を取り込む順序が堅実である。本特集の一貫した含意は、構造が効率的でもさばき方が粗ければ利益は漏れ、運用の設計こそが効率の上限を決めるということである。

\subsection{待ち行列の約束を守る}
運用における約束とは、計算量の少なさではなく、締め切りと同時進行の捌き方である。模型が軽くなっても、待ち行列の先頭で長い処理が居座れば、後続の応答はまとめて遅れる。利用者が感じる破約の多くは、演算器の不足ではなく、置き場所の競合から生じる。とくに生成推論では、入力を一気に読む段階と、一語ずつ吐き出す段階とで資源の使い方が異なり、両者を同じ列に無造作に並べると、一方の重さが他方の締め切りを壊す\autocite{splitwise}\autocite{sarathi}。

この競合をほどく鍵は、記憶を連続領域として抱え込まないことにある。鍵と値の履歴を細かな塊に分けて管理すれば、断片化で使えない隙間が生じにくくなり、空いた場所に次の要求を詰め続けられる\autocite{pagedattn}。重要なのは、単に詰め込み率が上がることではなく、優先度に応じた中断と退避に柔軟性が生まれることである。急ぐ応答のために譲る、長い入力のために一時的に外す、終わった分から回収する。そうした取り回しが利くからこそ、締め切り順の約束が運用として書ける。

長い入力の扱いは、その柔軟性を試す典型である。入力読みを一括で通すと、そのあいだ吐き出し側が止まり、体感遅延が跳ねる。これを小さな断片に区切って吐き出しの合間に相乗りさせれば、管の泡を埋めながら、吐き出し側の停滞を抑えられる\autocite{sarathi}。接頭辞の再利用も、待ちの短縮として読み直せる。共通の前置きを塊単位で使い回せば、同じ入力を繰り返し読む無駄が消え、列の進みが速くなる\autocite{sglang}\autocite{vllmprefix}\autocite{lmcache}。

最後に、受け入れの見方を分けておきたい。合成的な負荷で測った受け入れの良さは、あくまで測定条件つきの記録であり、そのまま運用の棄却挙動を統治しない。実運用では、締め切りを守れない要求は待たせるのではなく断る判断が要る。約束を守るとは、すべてを受け入れることではなく、受け入れられる範囲を見極め、溢れを早めに返すことである。

\subsection{受け渡しの経済学}
分離は、得意分野への特化で稼ぎ、受け渡しの費用を払う設計である。入力読みと吐き出しを別々の資源溜まりに置けば、それぞれの計算と記憶の性格に合わせた置き方ができる。ただし分離した瞬間に、履歴の受け渡しという新たな勘定が発生する。特化の利得から受け渡しの費用を差し引いた残りが、実質の利得である\autocite{splitwise}。

この勘定を理解するには、記憶を付随物ではなく一級の資源として扱う見方が要る。鍵と値の履歴を、計算の傍らに置かれた一時物ではなく、溜めて、集めて、運べる資源とみなせば、設計の選択肢が変わる。履歴中心の分離は、この資源観を前面に出したものであり、記憶の置き場所と運び方を運用の中心に据える\autocite{mooncake}。

外部化された再利用と退避の層は、同じ経済学を実行機関から切り離す。特定の実行機関に閉じない形で履歴の使い回しや退避や受け渡しを担えば、再利用の仕組みと退避の仕組みを組み合わせやすくなる\autocite{lmcache}。退避先の選び方も、経済計算の一部である。装置内の記憶に置ききれない履歴を、計算から少し離れた大きな記憶へ逃がせば、同時進行の余裕が生まれる。逃がす先を近くに置けば戻しは速いが置き場所の節約は小さく、遠くに置けば節約は大きいが戻しに時間がかかる。万能の正解はなく、置き場所ごとの出し入れの費用を見ながら決めるほかない。

運用の評価は、だから総量の伸びだけを見てはならない。分離で処理量が伸びても、受け渡しで食われた分を差し引かなければ、正味は分からない。受け渡しの経済学とは、特化の華やかさではなく、差し引きの地味さを直視する姿勢である。利得から費用を引いた残りが正である配置だけが、約束を守りながら残る。

\subsection{文書と実装の対応}
文書は生きており、ある固定構成での検証は、そのまま他の構成へ移らない。ある像とある設定とある同時実行の上限で確かめられた振る舞いは、その組み合わせの記録であって、別の像や別の上限での保証ではない\autocite{vllmrecipe}。手順書が生きて動く段階では、検証の記録は特定の時点の姿を留めたものにすぎない。固定された教科書として読むのではなく、変化する運用手帳として読む。どの組み合わせで確かめたのか、その確かめはどこまで移せるのか、その二点を欠かさず確認する読み方が要る。

設計の置き換えは、仕様と実行の対進化として現れる。従来の投機部品を捨て、別の方式の起草と検証へ切り替えるような変更は、模型側の取り替えだけで完結しない。実行側がその起草を受け止め、検証の流儀を合わせ、後方装置と歩調をそろえて初めて振る舞う\autocite{dspark}。核となる演算や通信や土台も、同じく組み合わせで読むべきものである。注意機構に寄り添う復号核、専門家並列の送受信を担う通信庫、低精度行列積を支える演算庫、そして頁化や疎や投機の経路を束ねる共有土台\autocite{flashmla}\autocite{flashinfer}\autocite{deepep}\autocite{deepgemm}。運用者の本当の関心は、技法の名前ではなく、どの組み合わせが振る舞うかである。文書に載る手順、模型側の来歴、実行側の受け止め、核と通信と土台の顔ぶれ、再利用と退避と受け渡しの置き方。これらがそろって初めて、ある配置は動く配置になる\autocite{vllmrecipe}\autocite{vllmv41}\autocite{vllmprefix}。

\end{multicols}

\begin{center}
\small
\begin{tabularx}{\textwidth}{@{}p{0.20\textwidth}p{0.34\textwidth}X@{}}
\toprule
運用の無駄と処方 & 仕組みの要点 & 払う代償・留保 \\
\midrule
記憶の断片化 & 塊単位の管理で再利用を細かく回す & 表設計の細部には立ち入らない \\
同じ前置きの再計算 & 接頭辞の木共有・外部層 & 粒度の違いは使い分け \\
段階の性質差 & 前半向きと後半向きに分ける & 受け渡し転送の費用 \\
分散保持 & 記憶を第一級資源として束ねる & 保持自体にも費用 \\
読み込みと生成の干渉 & 断片化して合間に滑り込ませる & 区切り方は条件依存 \\
復号の核 & 圧縮注意に寄せた実行 & 構造側と対の工夫 \\
専門家並列の通信 & 分散と結合のさばき & 計算と同等の重み \\
低精度の土台 & 細粒度の行列積 & 学習と運用の接着剤 \\
共有核 & 複数技法の同居 & 世代差の吸収 \\
生きた手順書 & 固定構成での確認 & 他構成への転用は慎む \\
\bottomrule
\end{tabularx}
\end{center}
\autocite{pagedattn,sglang,splitwise,mooncake,sarathi,flashmla}

\begin{claimboundary}[本節の読解上の境界]
設計の細部、条件付き評価、配置解析、費用表、手順の全体には踏み込まない。文書の対応は特定構成での確認に留まり、他構成への転用はしない。合成的な受け入れ確認は測定として扱い、実運用の拒否動作と切り分ける。
\end{claimboundary}

\section{「大きなモデルを毎回呼ぶ」必要はあるか}
\label{sec:routing}
\sectionkicker{ROUTING AND SPECIALIZATION}

\begin{multicols}{2}
\raggedcolumns
すべての問いをいきなり最大のモデルに投げるのではなく、安価で軽いモデルから順に試し、足りなければ強いモデルへ引き継ぐというのがカスケードの基本的な発想である\autocite{frugalgpt}。問いに対する採点器と、どのモデルへ振るかを決める振り分け器、そしてここで止めてよいかを判断する打ち切り判定を組み合わせ、精度を保ちながら全体の計算量や利用料を抑える。ここで見逃せないのは、料金体系がモデルごとに異なる点である。入力と出力のどちらに重く課金するかは一様ではなく、しかも料金の並び順自体が対象の課題によって変わる。つまりカスケードは単なる節約術ではなく、課題特性と料金構造を踏まえた割り当て問題として理解すべきものである。運用ではトークン単価ではなく課題完遂あたりの費用で捉える。ただし本稿では当時の料金体系に基づく報告として扱い、現在の価格にはそのまま当てはめない。最適化の定式化の全体や詳細な数表の再現には立ち入らず、仕組みと考え方に焦点を当てる。

\subsection{学習型の振り分けとその脆さ}
問いごとに強いモデルと弱いモデルを使い分ける学習型の振り分けでは、問いの難しさや性質を見極める予測器が要になる。代表的な手法では、人手の好みデータから勝敗予測を学び、強いモデルが必要な問いだけを強いモデルへ回す\autocite{routellm}。狙いは、簡単な問いを安価にさばき、難しい問いだけに高価な計算を充てることで、全体としての費用対効果を高めることにある。ただし好みの分布は時間とともに変わり、学習時の分布と運用時の分布がずれれば予測の当たり外れも変わる。判定を担うモデル自体の偏りや、評価用データの混入対策の限界も、結果の読み方に影響する。詳細な成果数表や構造の比較検討には立ち入らず、データと予測に基づく振り分けという設計思想と、その前提条件に焦点を当てる。したがって学習型振り分けは、万能の節約装置ではなく、分布の近さと見極め精度に依存する条件付きの仕組みとして捉えるべきである。

振り分けによる節約の話には、必ず留保条件を付けなければならない。好みデータで学んだ振り分け器は、分布が変わると脆さを見せることが指摘されている\autocite{fragility}。学習時と同じような問いが来る間は節約と精度の両立が成り立っても、話題や言い回し、難しさの分布がずれると、強いモデルに回すべき問いを弱いモデルで済ませたり、その逆が起きたりしやすくなる。この指摘は、先行する節約の報告を否定するものではなく、正しく使うための条件を付け足す役割を果たす。どのようなずれに弱いかは、要約として伝えられている範囲にとどまり、本文の詳細な条件分けには立ち入らない。一般化して語れる段階にはなく、条件の細部は未確認として扱うのが筋である。運用では定期的な見直しと、振り分けの当たり外れを監視する仕組みが前提になる。

\subsection{生成しない判断という特化策}
生成で長い文章を作る代わりに、状態を受け取って型付きの確率的な答えを返すという特化型の事例がある\autocite{sysone,sysoneconcepts}。並列に動作する標本抽出器と独自の学習手法の組み合わせで作られ、文字列の生成ではなく選択肢の選択や点数付けといった操作を直接行うと説明されている。学習手法の内部詳細は公開されておらず、外形的には並列性と出力の意味づけの範囲で理解するしかない。確信度は多くの予測を束ねた集計としての較正を意図したもので、個々の答えの正しさを保証するものではないと説明されている。入力は文章のみで英語中心であり、分類や振り分け、採点、抽出、分岐といった流れの中の判断部分に向くという位置づけであり、\textbf{自由な文章生成を置き換えるものではない}。較正の実質的な質や英語以外の精度は独立には確かめられておらず、手元の内容で試すことが前提とされている。契約の細部には立ち入らず、置き換えではなく使い分けの事例として読むべきである。

提供元は、この特化型について高速性と低廉さをうたっている\autocite{jevmodels}。短時間での応答や、従来の大型モデルに対する大幅な速度・費用の優位が、自社で組んだ業務処理の評価として示されているとされる。ただし測定は特定の手元の環境から行われ、比較対象の動かし方に独自の包みが挟まる点や、参照回答の偏りなども開示されており、数字を額面通りに受け取ることはできない。したがってこれらは確立した事実ではなく、\textbf{提供元による主張として受け止める}必要がある。独立した再現は未確認であり、速度や価格、処理量の主張はその留保付きで扱う。英語以外の精度は弱いことが明記され、手元の内容での確認が推奨されている。価格や処理上限は文書時点の記載であり、将来の変更によって変わり得る。

この特化型をめぐる外部の情報は、補強的な位置づけにとどまる。有志の覚え書きは一次文書と整合的な記録を残しているとされるが、仕様の典拠にはならない\autocite{jevwiki}。愛好者による再現の集積地も、類似の仕組みを試す手がかりを集めている段階とされ、査読を経た独立再現が確認されたわけでもない\autocite{awesomejev}。応答の型や確信度の導出など契約の細部は、照会先の資料に属する事柄であり本稿では立ち入らない\autocite{typesafeapi}。較正を確かめる公開手順も見当たらないというのが現状の記録である。以上を踏まえ、本稿ではこの事例を、定型の判断を安く速くさばくための現時点での特化策として位置づけ、生成一般の代替とは区別する。導入の可否は、自前の課題での試測と監視を前提に判断すべきである。

\subsection{閾値の調整が設計の核心である}
カスケードの実体は、大きな模型をいつ呼ぶかを決める判定装置である。採点器が候補回答の良し悪しを見積もり、振り分け器が次にどの模型へ回すかを選び、停止判定がここで打ち切るか続けるかを決める\autocite{frugalgpt}。採点器が甘ければ止めるべきでないところで止まり、厳しすぎれば止められるはずの問いまで上位へ送ってしまう。停止規則はその二つの誤りを最終的に引き受ける場所であり、閾値の設定が設計の核心になる。

選択肢は模型の並べ替えだけではない。問い方を変える工夫や、近似的な処理で済ませる手も同じ俎上に載る。問いを言い換えて安い模型でも解ける形に整える、少数の例を添えて出力を安定させる、検索や要約などの前処理で入力を短くする、といった手段は、いずれも精度と費用の交換関係を動かす。まず問いと前処理の側でできることを尽くし、それでも残る難問だけを上位へ送る順序が自然である。

費用の数え方も課題ごとに行う必要がある。入力と出力で単価が異なる料金体系では、長い生成を伴う課題と短い判定で済む課題とで有利な模型が入れ替わる\autocite{frugalgpt}。机上の単価比較で決めた構成は、運用に入った途端に崩れることが多い。閾値の調整は、緩めるか締めるかの方針選択である。緩めに止めれば費用は下がるが見逃しが増え、厳しく止めれば品質は守りやすいが上位呼び出しが増える。どちらが正しいかは課題の痛みによって変わる。一度決めた閾値を固定値として扱わないことである。問い合わせの分布が変われば、学習時の選好データとのずれが生じ、振り分け器の当たり外れの勘所もずれる\autocite{routellm}\autocite{fragility}。閾値は設計時に決めて終わりではなく、運用しながら直すものと考えるべきである。

そのため、カスケードには監視の循環を組み込んでおきたい。止めた問いと上げた問いの比率、上位へ送った問いのうち実際に品質が上がった割合、課題別の費用と誤り率を、版ごとに眺める。ずれが見えたら閾値を動かすか、そもそも振り分けの前提が崩れていないかを確認する。逆に監視なしで閾値だけを決めると、節約の仕組みが形骸化し、大きな模型を毎回呼ぶ構成と実質的に変わらなくなる。閾値の運用こそがカスケードの本体である。

\subsection{確信度との付き合い方}
確信度付きの出力を扱う際の出発点は、その意味を取り違えないことである。ここでいう較正は、多くの予測を束ねたときに、確信度に見合った頻度で当たりが出るという集計的な性質を指す。個々の回答について、この確信度だからこの回答は正しいという保証を与えるものではない。この点は提供側の説明自体が明示している注意であり、利用者側が勝手に読み替えてはならない\autocite{sysoneconcepts}。確信度は判断の材料にはなるが、判断そのものではない。

実務的な使い方は、分解して上げていく形に寄せるのが無理がない。大きな問いを小さな判定に割り、分類や振り分けや採点や抜き出しや分岐といった原子的な問いにして、組み立ては確定的な処理側で行い、確信度が低いものだけを大きな模型へ上げる。この分解と上げ方の指針は、提供側の概念説明が示す想定用途そのものである\autocite{sysoneconcepts}。確信度はその境界を引くための物差しであり、境界の引き方は課題の責任の重さで変える。型があるから正しいのではなく、型があるから運用の手当てがしやすい。その違いを見失うと、型への安心感が精度への過信にすり替わる。

日本語を含む非英語の扱いには明確な弱さが残る。英語中心で、非英語では精度が落ちると説明されており、自分の内容で試すよう注意が添えられている\autocite{jevmodels}。多言語の注意書きを軽く見て英語での感触のまま閾値を決めると、日本語の運用で見逃しが増える。日本語の読者を想定する本調査の文脈では、この点は譲れない。日本語の実データで小さく試し、言語別の誤り率と確信度の対応を確かめてから基準を決める必要がある。較正の振る舞いについても、公開された較正手順が見当たらず、独立した裏づけは見つかっていないという限界が残る。確信度の数字が工程表に載っていても、その数字の作られ方と確かめ方が開かれていない以上、数字を額面で受け取って自動判定に直結させるべきではない。

試し方の規律も欠かせない。まず範囲を絞り、代表的な問いを集めて、確信度ごとの的中率を自分の分布で描く。全体の的中率だけを見ず、確信度の帯ごとに分けて眺める。運用で流れる比率に合わせて標本を組まないと、試験で決めた閾値が現場で機能しない。別名で版が切り替わると背後の挙動が変わり得るという指摘もあり、応答に含まれる版表示を記録に残して閾値の履歴と突き合わせる工夫が有効である\autocite{jevmodels}。小さく試して小さく直す循環を回せるかどうかが、確信度付きの部品を使いこなせるかどうかの分かれ目になる。

\subsection{前段のさばき役としての居場所}
全体像のなかでこの種の特化型の居場所を考えると、前段のさばき役がしっくりくる。すべての問いを置き換えるものとしてではなく、振り分けの前段で易しい問いをさばき、難しい問いを後段へ渡す役として置く発想である。置き換えかさばき役かという枠組みの違いは、評価の問い方を変える。置き換えの枠組みでは平均精度の優劣が問われるが、さばき役の枠組みでは、さばいた範囲での誤り率と、後段へ渡した範囲での改善と、全体の費用と手間の合計が問われる。特化型を単体の賢さで値踏みすると見誤る。全体の流れのなかでどれだけ手間を減らしたかで見るべきである。

導入の判断は、自分の分布での試験に委ねるほかない。公開の説明は提供側による報告として受け止め、独立した再現が未了であることを前提に扱う。関連する再現の試みは集約の手がかりとしてのみ位置づけられ、一件ずつの精査や査読付きの再現としては確立していない\autocite{awesomejev}。照合のための覚書的な記録は参考にはなるが、仕様の典拠として引用すべきではなく、内容の変わりやすさにも注意が要る\autocite{jevwiki}。分布のずれのもとでは振り分けの利得が条件付きになるという要約段階の注意も、この試験優先の姿勢を後押しする\autocite{fragility}。

費用の勘定には監視の手間も含める。閾値の見直し、標本の作り直し、版切り替わりの確認、言語別の品質点検は、いずれも運用費である。これを削って部品単価だけを数えると、安く見えて高くつく構成になる。大きな模型を毎回呼ぶ必要があるかという問いへの答えは、呼び方を工夫すれば減らせるという一般論に尽きない。どこで止め、どこで上げ、どこで見直すかを決めた運用だけが、その答えを持ち帰れる\autocite{frugalgpt}\autocite{sysone}\autocite{sysoneconcepts}。

\end{multicols}

\begin{center}
\small
\begin{tabularx}{\textwidth}{@{}p{0.20\textwidth}p{0.28\textwidth}p{0.28\textwidth}X@{}}
\toprule
運用の問い & 止める判断 & 上げる判断 & 見直しの合図 \\
\midrule
閾値の置き方 & 緩めれば費用減と見逃し増 & 厳しめなら品質保持と費用増 & 分布ずれで勘所がずれる \\
確信度の読み & 帯ごとの的中率で描く & 低い帯だけ上位へ渡す & 言語別・版別の対応を確認 \\
問いと前処理 & 言い換えと例示と短縮を尽くす & 残る難問だけ上位へ & 前処理の改善が先 \\
費用の数え方 & 入出力の量と回数を含む & 課題完遂あたりで比べる & 単価表の一律優劣は捨てる \\
監視の循環 & 止めた割合と改善率を眺める & 前提崩壊の有無を確かめる & 形骸化の前に閾値を動かす \\
\bottomrule
\end{tabularx}
\end{center}
\autocite{frugalgpt,routellm,fragility,sysoneconcepts,jevmodels}

\begin{center}
\small
\begin{tabularx}{\textwidth}{@{}p{0.22\textwidth}p{0.34\textwidth}X@{}}
\toprule
呼ばないための設計 & 仕組みの要点 & 条件と留保 \\
\midrule
段階的な引き継ぎ & 安価な順に試し不足分だけ上位へ & 当時の料金体系の報告 \\
学習型の振り分け & 好みデータから勝敗予測を学習 & 分布の古びに依存 \\
振り分けの脆さ & 分布ずれで当たり外れが変わる & 要約範囲、一般化しない \\
型付きの特化判断 & 並列標本と確信度つき選択 & 生成の代替ではない \\
速度・費用の主張 & 自社評価での優位の提示 & 提供元主張、再現未確認 \\
確信度の扱い & 集計としての較正の意図 & 個の正しさの保証ではない \\
外部の覚え書き & 一致の記録 & 仕様の典拠にしない \\
再現の集積地 & 手がかりの収集段階 & 確立した裏付けなし \\
較正の手順 & 公開手順が見当たらない & 読み替えに慎重 \\
\bottomrule
\end{tabularx}
\end{center}
\autocite{frugalgpt,routellm,fragility,sysone,jevmodels}

\begin{claimboundary}[本節の読解上の境界]
料金は文書時点の記載であり現在価格に当てはめない。脆さの条件細部は未確認として一般化しない。特化型は現時点での特化策であり、生成一般の代替として語らない。独立再現と較正手順の確認までは代替として扱わない。
\end{claimboundary}

\section{2026年の実装点}
\label{sec:capstone}
\sectionkicker{2026 CAPSTONES}

\begin{multicols}{2}
\raggedcolumns
この節は四つの実装点を順位付けしない。優劣表を作ることが目的ではなく、各実装がどの層で計算と記憶と通信を削っているかを機構の積み重ねとして読み解くことが目的である。比較軸は総パラメータと作動パラメータ、注意機構と記憶保持の設計、数値精度、復号の直列性、サービングと実行環境、手元での実行可能性、文書化された代償、そして証拠の質である。四者の証拠の深さはそろっていない。本文まで構造的に読み込めた報告もあれば、要旨の範囲に限定される報告もある。ベンダー側の比率や速度表現は測定条件付きの主張として扱い、独立検証のない数値は事実として引用しない。実装リポジトリは存在証明の水準にとどめ、コード監査や独自の再実行は行っていない。手順書類は更新でずれる生きた文書であることを前提にする\autocite{v41report,qwen38next,glm5report}.

\subsection{大規模疎活性化と記憶圧縮を重ねた設計}
一つ目の実装点は、数百B級の総量に対して作動量を一桁B級前後に抑える疎活性化を土台に、注意と記憶の保持を同時に圧縮する設計である\autocite{v41report}。復号時と事前充填時で作動量を使い分ける条件付き計算、疎な注意の複数動作様式の共有、学習時から低ビット化された大域記憶、局所窓と再生機構による持続記憶の削減、条件付き記憶の追加が重ねられている。チェックポイントの同一性表示は、因果符号化復号化器や第二世代の疎注意、条件付き記憶、推測復号の構成を含むものとして、模型面の文書で示される\autocite{v41card}。

数値精度の結び付きがサービング要件を決める。専門家部分と密部分と埋め込みで精度体系が異なり、低精度経路は対応する加速器と実行基盤を条件とする。復号の高速化は単一の推測復号方式に置き換えられ、従来の多トークン予測部は当該チェックポイントでは用いない設計選択となっている。配備手順は集約配置と分離配置の双方を試験範囲として文書化する\autocite{dynamo}。推論努力量の数値対応や特定加速器での検証拒否時の代替動作まで含むが、手順書は更新される生きた文書であり、特定構成での検証は他構成や他画像に転用できない\autocite{vllmrecipe}。受容率の合成測定は測定値であり、実サービングの棄却動作とは区別する。実行文書の対応は、模型符号の文書として伝えられる\autocite{vllmv41}。

\subsection{残差と最適化と疎注意を共同設計した構成}
二つ目の実装点は、百B級強の総量に対して作動量を一桁B級前後に抑えつつ、ホスト側に置かれる表を組み合わせる構成である\autocite{qwen38next}。報告本文の範囲で読み取れる要点は、候補変更を損失と指標、学習と事前充填と復号の費用、安定性への影響の三軸で評価する姿勢、分岐と素子ごとのゲートを持つ残差変形、表の配置と語彙配分の検討、対象パラメータごとの最適化手法の割り当てである。損失と下流精度が乖離し、表の拡大が損失を単調に下げても精度は飽和するという観察が共同設計の主張を支える。設定表は層様式の比率や単一の予測層の存在を特定するが、自己報告の性能表は条件検証を待つ主張として隔離する。公式解説は刷新の枠組みを示す告知水準であり、実質は報告本文の裏付けに依拠する\autocite{qwenblog}。実装集積所は派生版の区別を結び付け、運用版とプレビュー版の混同を防ぐことが引用上の要点である\autocite{qwen38repo}。注意機構の設定や評価値は、公開の読み物としても示される\autocite{qwenreadme}。

実行側は、特定加速器での配置報告と地域計算環境への移植の存在で裏付けられる。加速器上の報告は接頭辞一致率を条件とする配置の存在証明としてのみ扱い、倍率表現はキャッシュ局所性の測定でもある点に注意して事実化しない\autocite{nvblogqwen}。報告面の取得は行っておらず、数値の独立検証やハーネス分解も未了である。地域計算環境への移植は実装存在の権威であり、変換器や低ランク結合を含む文グラフ、多数専門家からの上位数選択、予算付き疎注意、ホスト側の表処理などが実装記述として確認される\autocite{qwenpr}。検証方法や既知の実装上の隙間も報告者の記述として透明化されているが、独自の再構築や測定は行っておらず、移植の正しさは報告ベースの検証に依拠する。

\subsection{線形注意系の実行可能性}
三つ目の実装点は、注意の二次計算を線形系の記憶更新に置き換える系列に属し、互換性と運用面の配慮を前面に出す設計である。実装集積所は核計算と実行統合経路とチェックポイントの公開を一体として示し、報告を超える配備権威の役割を担う\autocite{kimigithub}。ただし集積所の複製や監査は行っておらず、独立再現は未実行である。核実装と実行経路の組み合わせは、報告を超える実行可能性の証明として配置水準で評価するが\autocite{kda}、核の数値計算やサービング性能は未消費であり事実として引用しない。名称衝突の可能性がある別系列の機構とは切り分け、混同しない。運用者の受容報告は補助的な位置付けにとどめ、技術的事実や経済的事実の根拠には用いない。

\subsection{要旨範囲に限定される四つ目の実装点}
四つ目の実装点は証拠の扱いが異なることを最初に明示する。技術報告については要旨のみを消費し、本文の採用詳細や強化学習基盤の設計や評価表は未消費である\autocite{glm5report}。したがって本文で裏付けられた言語では記述しない。数百B級の総量に対して作動量を数十B級前後に抑える基礎模型として、他所発の疎注意機構の採用と言語化されるが、その意義付けは要旨範囲の採用事実にとどまる。混合深度として、チェックポイントの同一性表示\autocite{glm53card}、開発者文書の相対比率\autocite{zaidocs}、実装集積所の共有機構と予測層と言及される配置表を組み合わせる\autocite{glm5github}。開発者文書の計算量や記憶量の相対比率は、基準を名指したベンダー相対の測定として条件付きで扱い、測定条件は未確認として隔離する。集積所の数値は集積所報告を超えて裏付けない。未解決のプール名称は事実に用いない。模型面の更新表現や順位表現は二次的として事実化しない。起動時の出来事は告知水準として位置付け、抜粋要約の範囲を超えない\autocite{dsadebut}。

手元実行は量子化物の配布と記憶適合指針の存在が要点であり\autocite{unslothglm}、量子化品質表や速度主張は提供者測定として独立確認を待つ。横断すると、復号高速化の置換判断、手元実行の現実性、効率から価格への波及の三点が代償と証拠の質を示す。推測復号の方式置換は意図的な設計選択であり、受容に関する数値は特定負荷での測定に限定される\autocite{dspark}。生きた手順書の変動、合成受容と実棄却の区別、条件付き比率の隔離、再現の未了を前提にすれば、四者を機構の選択の違いとして理解できる。順位ではなく、どの隘路を何で削り何を代償にしたかが結論である。

\subsection{作動量の設計思想の違い}
総量と作動量の分離は、この年の実装点に共通する関心であるが、その切り分け方は各設計で異なる。DeepSeek系のFlashは、符号化時と生成時で働く量を変える条件付き計算を中核に据える\autocite{v41report}\autocite{v41card}\autocite{vllmrecipe}。常に全体を動かすのではなく、場面ごとに必要な部分だけを作動させる思想であり、報告書本文と配備用手順書の双方で構造として確認できる。この選択は、常駐と通信の設計に直結する。動かさない部分の重みや記憶をどこに置くか、分離配置における鍵値受け渡しをどう担うか、専門家並列の振り分けをどの実装経路でさばくかといった配備条件が、作動量の思想と一体で記述される\autocite{dynamo}\autocite{vllmrecipe}\autocite{vllmv41}。作動量を絞ることは単なる節約ではなく、置き場所と運び方の再設計を要求する。

Qwen系の次世代Flashは、残差経路のゲート制御による絞り込みと、表参照による補完を組み合わせることで、総量と作動量を分ける。重みとして保持する知識と、表として引く知識を別系統に置き、計算で解く部分と参照で済ませる部分を分担させる共同設計が本文で主題化されている\autocite{qwen38next}\autocite{qwenreadme}\autocite{qwenblog}。表は主記憶とは異なる場所に常駐し、必要に応じて引かれるため、加速器側の作動量は抑えられる一方で、ホスト側の表保持と変換時の作業記憶が実装上の関心となる。この振る舞いは移植報告における検証記述からも読み取れ、表の行照合や分割前処理の同一性確認が移植正当性の柱に置かれている\autocite{qwenpr}\autocite{qwen38repo}。計算を減らす代わりに置き場所を増やすのか、置き場所を集約する代わりに計算を許容するのかという分岐が、ここでは明確に表の有無として現れる。

Kimi系の線形系は、作動量の重心を注意の選択そのものから外す方向を選ぶ。二次的に増える注意行列を保持・移動させるのではなく、漸化的に更新される状態に記憶を畳み込むことで、層ごとの選択的作動という発想を置き換える。公開リポジトリがカーネル実装と統合経路、チェックポイントの公開を一体で示すことは、この思想が実装可能性と切り離せないことを物語る\autocite{kimigithub}\autocite{kda}。常駐するのは注意の断片ではなく圧縮された状態であり、通信するのは点数行列ではなく状態更新に必要な最小限の情報となる。ただしカーネル実装の数値や性能は未消費であり、思想としての方向性と配備可能性の存在は語れても、運用像の断定には進まない。この留保は欠落ではなく、証拠の深さに応じた言葉の選び方である。

GLM系のFlashに関する相対比は、隔離されたベンダー主張として扱い、事実としては引用しない。注意計算量や鍵値量の削減を示す相対表示は、測定条件が未確定の文書面での記述に留まるため\autocite{zaidocs}\autocite{glm5github}、本稿では削減幅の多寡を語らず、相対比が単独では設計思想の比較材料にならないことのみを確認する。推論努力量の段階制御も、作動量とは別軸の設計である。手順書が示す努力量の概念は、復号の振る舞いや推論の深さを運用時に切り替えるためのもので、重み側の条件付き計算と混同できない\autocite{vllmrecipe}\autocite{dspark}。いずれの設計も優劣で並べず、分離の仕方とその代償の置き方の違いとして読むことが、本節の契約である。

\subsection{注意と記憶の選択の違い}
疎な索引型、予算付き疎型、線形記憶型、他所発機構の抽象採用型という四つの注意選択は、何を残し何を動かすかという点で対照的である。DeepSeek系のFlashが採用する疎な索引型は、大域の鍵値を残したまま索引機構で選ぶ方式である\autocite{v41report}\autocite{vllmrecipe}\autocite{vllmv41}。残すのは全体であり、動かすのは選ばれた分だけである。ここに窓付き注意と再生機構が加わる。近傍窓で局所性を担保しつつ、再生機構で永続的な占有を抑えることで、長期文脈における常駐の増大を構造的に抑える。圧縮を担う層を限定する設計も、同様の発想の延長にある。すべてを一様に圧縮するのではなく、効果の見込める箇所に限定することで、保持と移動の総量を抑える。

Qwen系の予算付き疎型は、予算という物差しで密と疎を切り分ける。予算内に収まる範囲では密計算と等価であることが移植報告の検証で確かめられ、予算を超えた範囲で選択的注意に移る\autocite{qwen38next}\autocite{qwenpr}\autocite{nvblogqwen}。残すのは層全体の表現力であり、動かすのは予算で区切られた有効分である。この方式の帰結は、接頭辞一致率への依存として現れる。ベンダー側の高速化表示が接頭辞一致を条件とするのは、注意機構の測定であると同時に局所性の測定でもあるからであり、本文脈では数値の多寡を引用せず、条件付きであるという読み方だけを残す\autocite{nvblogqwen}\autocite{qwenblog}。予算付きであることは、無条件の削減ではなく、文脈の再利用のされ方に依存する削減であることを意味する。

Kimi系の線形記憶型は、注意を残さない選択である。鍵値を残して選び直すのではなく、状態に畳み込んで逐次更新することで、記憶の保持形式そのものを変える。第三者実装としてのカーネル実装と統合経路の存在は、この置換が文書上の提案に留まらず実行可能な形を持つことを示す\autocite{kimigithub}\autocite{kda}。ただし数値的性質や提供性能は未消費であり、保持形式の転換という設計事実と、運用品質の確定という事実を区別する。GLM系における他所発機構の採用は、移植可能性の事実として読む。DeepSeek系に由来する機構が別系統の旗艦に取り込まれたことは、機構の可搬性を示す独立した傍証であり、報告書要旨の範囲で確認できる\autocite{glm5report}\autocite{glm5github}。保持と移動の内訳に踏み込んだ記述はここでは行わず、採用という選択の存在とその意義に限定する。四者は序列ではなく、残すものと動かすものの配分の違いとして並置される。

\subsection{証拠の深さの違いと読み方}
本文消費、要旨限定、告知水準、移植報告という四つの深さは、許される言葉の範囲を変える。本文まで消費された資料は、構造と実装条件を語ることができる。DeepSeek系の報告書本文における条件付き計算や疎な索引の構成、推論後学習の工程の骨格は、構造的範囲での読解を許す\autocite{v41report}。Qwen系の報告書本文における残差ゲートや表配置、最適化手法の割り当て、損失と精度の乖離という共同設計の主題も、同様に本文読解としての言葉を許す\autocite{qwen38next}。配備手順書を全文消費した場合も、検証済み配置や推論努力量の意味づけといった運用条件を事実として語れる\autocite{vllmrecipe}\autocite{dynamo}\autocite{vllmv41}。ただし、いずれの本文消費においても、個別表の数値や詳細な除去実験の内訳、評価基盤に結びついた検証を欠く比較表示は引用しない。構造を語れることと、数値を事実として語れることは別の事柄であり、本稿では前者に留める。GLM系については報告書要旨のみを消費した範囲で論じ、本文の詳細や学習基盤の設計、表の内訳に踏み込んだ断定は行わない\autocite{glm5report}。

告知水準の資料は、方向と構想の表明として読む。公式発信における刷新点の枠付けや、前任者との対比による向上の言い回しは、ベンダー側の枠付けであって測定事実ではない\autocite{qwenblog}\autocite{nvblogqwen}\autocite{dsadebut}。加速器上の高速化表示も、接頭辞一致などの条件と一体で読むべき局所性依存の表示であり、条件を離れた効率の事実としては引用しない。移植報告水準の資料は、実装存在の証拠として読む。共同体移植における変換器や疎注意の実装、表参照のホスト側処理の記述は、実行可能性の存在を示す\autocite{qwenpr}\autocite{qwen38repo}。量子化文書における形式一覧や適合指針の記述も、同様に存在と手順の証拠である\autocite{unslothglm}\autocite{kimigithub}\autocite{kda}。ただし移植正当性は報告者の検証条件に結びついた検証に依拠し、独立再構築や独立測定を経ていないため、正確性の断定には進まない。量子化品質の表示は測定者限定の表示であり、高速化表示はベンダー測定であって確認を要することが境界として残る。

未解決の固有名は事実として用いない。GLM系リポジトリに現れる未解決の共有資源名は、同一性が定まらないためいかなる事実にも引用しない\autocite{glm5github}。本稿でいうカーネル実装は線形系に属する実装を指し、同名の別分野の機構との混同をあらかじめ排除する\autocite{kda}\autocite{kimigithub}。手順書の変動も読み方の一部である。配備手順は更新を前提とする生文書であり、固定像での検証は他像や他予定上限に転用できない\autocite{vllmrecipe}\autocite{dynamo}。合成受容と実棄却の区別も同様である。処理能力測定様式における合成受容は測定に留まり、実提供を支配するのは実棄却の挙動であるため、受容数の多寡を運用品質として引用しない\autocite{vllmrecipe}\autocite{dspark}。深さに応じた言葉を守ることが、順位付けを排して実装点を読むための条件である。

\end{multicols}

\begin{center}
\small
\begin{tabularx}{\textwidth}{@{}p{0.14\textwidth}p{0.24\textwidth}p{0.24\textwidth}X@{}}
\toprule
比較の軸 & 一つ目の設計 & 二つ目の設計 & 三つ目・四つ目 \\
\midrule
総量と作動量 & 数百B級と一桁B級前後 & 百B級強と一桁B級前後 & 線形系・数百B級と数十B級前後 \\
注意と記憶 & 複数様式の共有と窓再生 & 予算付き疎注意と表 & 実行経路・要旨範囲の採用 \\
数値精度 & 部位別の精度体系 & 表と残差の共同設計 & 核の数値未消費、相対比率 \\
復号 & 単一推測方式への置換 & 移植報告の検証 & 要旨範囲の置換の限定 \\
運用と実行 & 集約・分離の手順書 & 加速器配置と移植 & 集積所の存在証明 \\
手元の可能性 & 手順書の対応 & 移植の存在 & 量子化物の配布 \\
証拠の深さ & 本文と手順の裏付け & 本文と移植の裏付け & 配置水準・要旨限定 \\
\bottomrule
\end{tabularx}
\end{center}
\autocite{v41report,qwen38next,kimigithub,glm5report}

\begin{claimboundary}[本節の読解上の境界]
順位付けは行わない。四つ目の技術報告は要旨のみの消費であり本文言語では記述しない。ベンダー比率・速度表現は条件付き主張として隔離する。手順書は生きた文書であり他構成に転用しない。未解決の名称・衝突は事実に用いない。
\end{claimboundary}

\section{数字をどう読むか}
\label{sec:measure}
\sectionkicker{MEASUREMENT LITERACY}

\begin{multicols}{2}
\raggedcolumns
推論の速さは一つの数字では語れない。まず応答の出だしまでの待ちと、生成が続く区間の一トークンあたりの間隔と、単位時間あたりの出力と、同時実行や到着率の負荷をかけたときの全体処理量を分けて捉える必要がある\autocite{mlperf,genaiperf}。文脈が長くなれば出だしの読み込みが重くなり、同時実行を束ねれば待ち行列の振る舞いが変わり、装置が変われば同じ模型でも数字が動く。推論に割く手間を増やせば出力が長くなり、単位時間あたりの量の見え方も変わるため、手間の指定を欠いた速さの比較は成り立たない。計測文書に即した道具立ては、これらの指標を負荷条件とともに記録することを求める\autocite{genaiperf}。同時実行数や要求の送り方、計測区間の切り方、対象装置の特定を欠いた毎秒あたりの量は、再現の手がかりにならない。文書上の注意書きが多い計測枠組みは、数値の出所としてではなく、比較が壊れる理由を教える教材として扱うのが適切である\autocite{llmperf}。後継として名指しされた存在についても、単独の製品としての姿は確認できておらず判断を保留し、断定しない。

\subsection{実課題の系譜を読む}
実課題の解決力を測ろうとすれば、現実の開発現場にある課題報告とそれを確かめる試験に行き着く\autocite{swebench}。ただし初期の集め方には壊れた課題や仕様の曖昧さが混ざり、試験を通っても直したことにならない場合や、何を直すべきか定まらない場合が残る。この妥当性の綻びを直視し、人手で選び直した精選版や、長期にわたる作業を想定して作り直した難所版が生まれたのは自然な流れである\autocite{swepro}。難所版では人手で確かめられた課題が複数の保管庫から集められ、公開用と商用と非公開の区分が設けられて、混入や漏えいへの備えが意識されている。ここで守るべき作法は、区分の異なる数字を混ぜて並べないことである。公開の場で見える成績と、別条件の区分の成績は、集め方も漏えいの備えも異なるため、素朴に比べられない。個々の足回りの結び付きの特定は留保されており、本文の詳細な手順の読み込みも残されている。要旨に即した設計の理解にとどめ、現物の再実行による裏取りに踏み込まない。したがって個別の数値は隔離し、ここでは妥当性をめぐる設計の系譜だけを学ぶ。見かけの順位を新たに作らないことが、この節の約束である。

端末に向かう代理作業の評価は、応答文の見栄えではなく、操作を終えた後の容器の状態で合否を決める点に特徴がある\autocite{tbench}。ただし同じ名を掲げていても、改訂を経れば課題集合や判定手順は別物になり得る。名は契約ではないというのが、版の違いが教える第一の教訓である。別の注意例として、評価の足回りの選び方で成績が動くことが記録されている。複数の足回りのうち良い方を採る流儀が残れば、数字は実力だけでなく足回りとの相性をも映す。標準化された正確さ側の枠組みが担うのもこの役割であり\autocite{lmharness}、枠組み名を欠いた正確さの言及は比較の土台を欠く。現物の保管庫や課題設定の細部までは立ち入らず、ここでは個別の数値を引用しない。

\subsection{長さと報告の構え}
長い文脈をうたう表示と、実際に使える長さは別の事柄である\autocite{ruler}。段階的に長さを上げながら、取り出しや飛び越しのある照合、まとめ上げ、問い応答を等級づける考え方は、このずれを可視化する。有効な長さを等級づけることは、能力の物差しであると同時に、記憶と移動の費用を問う効率の物差しでもある。本文の課題群や採点手順の細部の読み込みが残る以上、等級づけの内訳に踏み込んだ断定は避け、要旨に即した理解にとどめる。一方で広い評価を志す立場は、場面と指標と適用条件を掛け合わせ、多面的な報告の透明性を重んじる\autocite{helm}。販売側の表を検証する側の物差しもここにあり、条件の開示を求める姿勢が共通する。保持の費用と読み込みの重さを見据え、長い文脈の表示をそのまま実力と受け取らないことが肝要である。

数字を並べる前の作法を束ねる。第一に、利用場面の束ねを欠いた比較をしない。第二に、来歴の異なる区分を混ぜない。第三に、版と足回りと道具利用を添える。第四に、混入と課題完遂の費用を見る。出だしと生成中と全体を分け、版と足回りと道具と手間を添え、区分を混ぜず、混入と総費用を見る。この確認を欠いた数字は並べない。新たな順位付けを作らず、条件付きで読む姿勢だけを持ち帰る。

\end{multicols}

\begin{center}
\small
\begin{tabularx}{\textwidth}{@{}p{0.24\textwidth}p{0.34\textwidth}X@{}}
\toprule
読みの作法 & 確認すること & 混ぜてはいけないもの \\
\midrule
局面を分ける & 出だし・間隔・全体処理量 & 待ちと滑らかさと捌き量の一括 \\
場面を束ねる & 到着・遅延上限・装置の同一性 & 条件欠きの毎秒量 \\
区分を分ける & 公開・商用・非公開の来歴 & 区分違いの成績の並置 \\
版を添える & 改訂・足回り・道具・手間 & 名だけの言及 \\
混入に備える & 集め方の備えと漏えい対策 & 暗記の映し込み \\
総費用で見る & 速さと正確さと手間 & 片手落ちの比較 \\
長さを疑う & 公称上限と有効長のずれ & 表示のままの実力視 \\
報告を開く & 場面・指標・適用の開示 & 代表値への畳み込み \\
\bottomrule
\end{tabularx}
\end{center}
\autocite{swebench,swepro,tbench,ruler,mlperf}

\begin{claimboundary}[本節の読解上の境界]
個別の数値は隔離し、設計の系譜と作法だけを学ぶ。現物の複製や再実行には踏み込まない。後継の単独姿は未確認として保留する。新たな順位付けは作らない。
\end{claimboundary}

\section{結び──層として読む効率}
\label{sec:synthesis}
\sectionkicker{FINAL SYNTHESIS}

\begin{multicols}{2}
\raggedcolumns
本巻がたどったのは、効率化が単一の物差しから層の重なりへと変わる過程である。規模拡大の経験則を起点に、条件付き計算への分岐、注意・記憶の共有化と疎化・線形化、重みの低ビット化と局所実行、逐次復号の先読みと推論予算の管理、運用の束ね・再利用・分離、安価計算への振り分けと特化型の使い分けが、各層の律速に応じて積み重なる。容量と計算の分離が全層を貫く転換であり、総量の競争から稼働量・常駐量・保持量の設計へ関心が移った。

並行する競合関係も見えた。密な規模拡大と疎な条件付き計算、厳密な入出力最適化と計算・保存を変える構造、学習後量子化と原生低ビット、ドラフト検証による待ち削減と推論予算の節約、集中運用と分離運用、汎用生成と型付き特化判断。いずれも万能ではなく、作業負荷・文脈長・同時実行・装置条件で損得が変わる。\textbf{どの隘路を、どの条件のもとで、何で削ったか}。この問いを手放さないことが、効率の議論を曇らせない条件である。

未解決の系譜は残る。外部記憶表を持つ設計と条件付き記憶提案の導出関係、相手陣営の由来文書の公表時期、呼称の衝突、特化型の独立再現と較正手順、後継製品の単独姿、節単位の消費が残る方法論、頂点実装の図別ピンは、いずれも未解決のまま残す。境界の明示は弱さではなく、検証の順序を守る手続きである。数字を読む作法を身につけ、条件付きの主張は条件付きのまま受け止めること。それが本巻の持ち帰りである。

層の読み方を実践に移すとき、三つの順序が助けになる。第一に、削減対象の名指しである。学習か推論か、前処理か復号か、演算か容量か帯域か通信か、遅延か処理能力か、総量か有効量か、重みかキャッシュか、トークン単価か仕事単価か。対象を名指ししないまま手法を選ぶと、効くはずの条件と外れる。第二に、条件の束ねである。到着のさせ方や遅延の上限、同時実行や文脈長、装置の顔ぶれ、版と足回りと道具と手間をそろえて初めて、二つの数字は並べられる。第三に、深さに応じた言葉である。本文まで読んだこと、要旨だけ読んだこと、告知だけ聞いたこと、移植の存在だけ確かめたことを区別し、それぞれに許される言葉で語る。この三つを守れば、効率の議論は曇らない。
\end{multicols}

\begin{claimboundary}[本巻の読解上の境界]
本巻は順位付けも推奨も行わない。概要水準の主張は概要水準のまま扱い、提供元測定は隔離する。未解決事項を閉じないことが、本巻の読みの作法である。
\end{claimboundary}

\Needspace{15\baselineskip}
\section*{Source-backed Technical Notes / 出典に裏付けられた技術要点}
\small
各節の技術的主張と読解上の境界を、一次資料の裏付けとともに表にまとめる。数値は条件付きの報告として扱い、横断の優劣には用いない。

\subsection*{第1節──隘路の契約}
\begin{tabularx}{\textwidth}{@{}>{\raggedright\arraybackslash}p{0.17\textwidth}>{\raggedright\arraybackslash}p{0.38\textwidth}>{\raggedright\arraybackslash}X@{}}
\toprule
論点 & 一次資料で確認する技術要点 & 読解上の境界 \\
\midrule
規模則 & 損失が規模・データ・計算に滑らかに従う関係の形 & 係数値・実験条件には立ち入らない \\
配分の見直し & 過小学習の指摘と均等配分の処方箋 & 目安を普遍定数として扱わない \\
段階の分離 & 反復単位の継続的束ねと前処理・復号の分離 & 処理細部・速度幅には踏み込まない \\
\bottomrule
\end{tabularx}
\autocite{kaplan2020,chinchilla,orca}

\subsection*{第2節──条件付き計算の系譜}
\begin{tabularx}{\textwidth}{@{}>{\raggedright\arraybackslash}p{0.17\textwidth}>{\raggedright\arraybackslash}p{0.38\textwidth}>{\raggedright\arraybackslash}X@{}}
\toprule
論点 & 一次資料で確認する技術要点 & 読解上の境界 \\
\midrule
疎ゲートの起点 & トークンごとの専門家選択と補助損失 & 数式・実験表の細部は範囲外 \\
並列と分割 & 分散配置と容量係数・管理可能性 &  bench横断の数値比較はしない \\
振り分けの分類 & 一人選択・側選択・粒度・共有 & 方式の優劣の断定はしない \\
稼働量の経済学 & 総量と稼働・常駐の分離 & 常駐・通信の費用を直視する \\
学習の土台 & 低精度の大規模学習と予測目的 & 速度の定量主張はしない \\
\bottomrule
\end{tabularx}
\autocite{shazeer2017,gshard,switcht,deepseekmoe,deepseekv2,deepseekv3}

\subsection*{第3節──注意と記憶}
\begin{tabularx}{\textwidth}{@{}>{\raggedright\arraybackslash}p{0.17\textwidth}>{\raggedright\arraybackslash}p{0.38\textwidth}>{\raggedright\arraybackslash}X@{}}
\toprule
論点 & 一次資料で確認する技術要点 & 読解上の境界 \\
\midrule
共有化 & 単一頭・束共有と移行可能性 & 代償の数値表は範囲外 \\
厳密な高速化 & 分割と逐次正規化の両立 & 証明・速度表は範囲外 \\
疎化と局所 & 索引選択と窓基準の対比 & 名称の混同を避ける \\
線形系譜 & 門付き差分と交互配置 & 再現は未確定 \\
保持と再利用 & 除去・共有・条件付き記憶・深さ & 数値比較はしない \\
\bottomrule
\end{tabularx}
\autocite{mqa,gqa,flashattn,nsa,kimilinear,engram}
\normalsize

\Needspace{17\baselineskip}
\subsection*{第4節──精度と局所実行}
\begin{tabularx}{\textwidth}{@{}>{\raggedright\arraybackslash}p{0.17\textwidth}>{\raggedright\arraybackslash}p{0.38\textwidth}>{\raggedright\arraybackslash}X@{}}
\toprule
論点 & 一次資料で確認する技術要点 & 読解上の境界 \\
\midrule
重み限定の圧縮 & 外れ値分離と再構成誤差の抑制 & 活性値同時・学習時と混同しない \\
同時量子化 & 平滑化と重要度保護 & 導出細部は範囲外 \\
器の理解 & 表現・コンテナ形式としての共有 & 量子化手法と呼ばない \\
低ビット浮動小数点 & 混合精度と微小スケーリング & 査読・断定はしない \\
適応と原生 & 付加器と一ビット方向 & 配備の証拠として扱わない \\
手元実行 & 不均質配置と移植報告 & 配置依存の観察に留める \\
\bottomrule
\end{tabularx}
\autocite{llmint8,gptq,qlora,gguf,fp8formats,bitnet}
\normalsize

\subsection*{第5節──復号と予算}
\begin{tabularx}{\textwidth}{@{}>{\raggedright\arraybackslash}p{0.17\textwidth}>{\raggedright\arraybackslash}p{0.38\textwidth}>{\raggedright\arraybackslash}X@{}}
\toprule
論点 & 一次資料で確認する技術要点 & 読解上の境界 \\
\midrule
投機的検証 & 草案と一括検証、分布保存 & 証明・速度表は範囲外 \\
先読みの質 & 特徴量予測と動的木、学習の改善と再利用 & 数値の一般化は留保 \\
予算の足し算 & 検証器探索と予算強制 & 訓練費と推論費を分ける \\
予算の引き算 & 上限制約の学習と過思考の知見 & 対象外への一般化はしない \\
\bottomrule
\end{tabularx}
\autocite{specdecode,eagle,mtp,snell2024,thinkprune,overthink}
\normalsize

\subsection*{第6節──運用}
\begin{tabularx}{\textwidth}{@{}>{\raggedright\arraybackslash}p{0.17\textwidth}>{\raggedright\arraybackslash}p{0.38\textwidth}>{\raggedright\arraybackslash}X@{}}
\toprule
論点 & 一次資料で確認する技術要点 & 読解上の境界 \\
\midrule
記憶管理 & 塊単位と連続的束ね & 表設計の細部は範囲外 \\
使い回し & 木共有と外部層 & 粒度の違いは使い分け \\
分離 & 段階別の資源と受け渡し & 転送費用を差し引いて評価 \\
噛み合わせ & 断片化と釣り合い & 区切り方は条件依存 \\
基盤 & 核・通信・低精度の三層 & 数値は引用しない \\
\bottomrule
\end{tabularx}
\autocite{pagedattn,sglang,splitwise,mooncake,sarathi}
\normalsize

\Needspace{15\baselineskip}
\subsection*{第7節──振り分けと特化}
\begin{tabularx}{\textwidth}{@{}>{\raggedright\arraybackslash}p{0.17\textwidth}>{\raggedright\arraybackslash}p{0.38\textwidth}>{\raggedright\arraybackslash}X@{}}
\toprule
論点 & 一次資料で確認する技術要点 & 読解上の境界 \\
\midrule
引き継ぎ & 安価な順の試行と打ち切り & 当時の料金の報告 \\
学習型振り分け & 好みデータの勝敗予測 & 分布依存の条件付き \\
脆さ & 分布ずれの当たり外れ & 一般化しない \\
特化判断 & 並列標本と確信度 & 生成の代替ではない \\
主張の扱い & 速度・費用の優位の提示 & 提供元主張、再現未確認 \\
\bottomrule
\end{tabularx}
\autocite{frugalgpt,routellm,fragility,sysone,jevmodels}
\normalsize

\subsection*{第8節──実装点の照合}
\begin{tabularx}{\textwidth}{@{}>{\raggedright\arraybackslash}p{0.17\textwidth}>{\raggedright\arraybackslash}p{0.38\textwidth}>{\raggedright\arraybackslash}X@{}}
\toprule
論点 & 一次資料で確認する技術要点 & 読解上の境界 \\
\midrule
疎活性化の設計 & 総量と作動量の分離 & 手順書は生きた文書 \\
共同設計 & 残差・表・最適化の三軸 & 自己報告の性能表は隔離 \\
線形系の実行 & 集積所と実行経路 & 数値・性能は未消費 \\
要旨限定 & 要旨のみの消費 & 本文言語では記述しない \\
\bottomrule
\end{tabularx}
\autocite{v41report,qwen38next,kimigithub,glm5report}
\normalsize

\subsection*{第9節──測定の作法}
\begin{tabularx}{\textwidth}{@{}>{\raggedright\arraybackslash}p{0.17\textwidth}>{\raggedright\arraybackslash}p{0.38\textwidth}>{\raggedright\arraybackslash}X@{}}
\toprule
論点 & 一次資料で確認する技術要点 & 読解上の境界 \\
\midrule
速さの分離 & 出だし・間隔・全体処理量 & 条件欠きの比較はしない \\
実課題の系譜 & 報告と試験の手順、区分の分離 & 区分違いを混ぜない \\
版の特定 & 改訂・足回り・道具・手間 & 名だけの言及はしない \\
長さの等級 & 公称と有効のずれ & 内訳の断定は避ける \\
\bottomrule
\end{tabularx}
\autocite{swebench,swepro,tbench,ruler,mlperf}
\normalsize

\Needspace{15\baselineskip}
\subsection*{用語集}
本巻で使う技術用語の定義を集めたものである。節をまたいで同じ言葉が同じ意味で使われていることを確かめるために置いた。
\begin{tabularx}{\textwidth}{@{}>{\raggedright\arraybackslash}p{0.22\textwidth}>{\raggedright\arraybackslash}X@{}}
\toprule
用語 & 本巻での意味 \\
\midrule
プリフィル & 入力全体を一気に読み込み内部状態を構築する前処理の段階。並列に処理しやすく演算が律速になりやすい \\
デコード & トークンを一つずつ逐次的に生み出す復号の段階。重みと中間記憶の読み出しが律速になりやすい \\
遅延 & 一つの応答が返るまでの待ち時間。対話の応答性に関わる \\
スループット & 単位時間あたりに処理できる要求量。同時接続の捌き量に関わる \\
総パラメータ & 模型が保持する重みの全体規模。保存容量と読み込み対象を決める \\
有効パラメータ & 一トークンの処理に実際に寄与する部分の規模。一トークンあたりの計算量の見通しを与える \\
KVキャッシュ & 逐次生成に伴って蓄積されるキーと値の中間記憶。文脈長と同時実行数に応じて膨らむ変動費 \\
条件付き計算 & トークンごとに一部の専門家だけを稼働させる計算様式。容量と計算を切り離す \\
専門家 & 条件付き計算で選択される部分模型。細分化と共有化の対象になる \\
振り分け & トークンや問いを専門家や模型に割り当てる操作。経路選択とも呼ぶ \\
受容率 & 草案が標的模型に受け入れられる割合。投機的復号の効率を決める物差し \\
ドラフト & 標的模型の検証を受けるために先に作られる草案。小さな模型や先読み層が担う \\
先読み & 未来の複数トークンを予測する学習と草案生成。補助課題と復号再利用の両面を持つ \\
予算強制 & 思考の打ち切りと延長で推論予算を操作する手法。短い思考の延ばしと長い思考の抑止を含む \\
過思考 & 長い推論が精度低下を招く現象。正答の反転を伴うことがある \\
塊管理 & 鍵と値の履歴を小さな塊に分けて管理する運用。断片化を抑え再利用の単位をそろえる \\
接頭辞再利用 & 共通の前置きの計算結果を使い回す運用。木共有と塊照合の二つの流儀がある \\
分離配置 & 前処理と復号を別の資源溜まりに置く運用。受け渡し転送の費用を伴う \\
カスケード & 安価な模型から順に試し不足分だけ上位へ引き継ぐ設計。打ち切り判定が核心になる \\
較正 & 確信度に見合った頻度で当たりが出る集計的性質。個の正しさの保証ではない \\
作動量 & 実際に計算に寄与する部分の量。総量との対比で効率を測る物差し \\
疎な注意 & 見るべき場所だけを選ぶ注意機構。索引選択と局所窓の二つの顔を持つ \\
線形注意 & 二次計算を漸化的な状態更新に置き換える注意系列。保持形式そのものを変える \\
量子化 & 数値のビット幅を減らして記憶と計算を軽くする操作。対象と時点の区別が要る \\
コンテナ形式 & テンソル配置と付随情報を定める器。量子化手法ではない \\
不均質実行 & 中央演算装置と画像演算装置に専門家を配置換えする運用。容量と計算を切り離す \\
出だしの待ち & 応答の最初の文字が出るまでの時間。TTFTとも呼ばれる \\
生成の間隔 & 生成が続く区間の一トークンあたりの時間。TPOTやITLとも呼ばれる \\
全体処理量 & 負荷をかけたときの単位時間あたりの出力。到着率と同時実行の条件と一体である \\
枠組み & 評価の道具立てと採点手順。枠組み名は数字の一部をなす \\
版 & 課題集合や判定手順の改訂。名が同じでも版が違えば別の物差しである \\
\bottomrule
\end{tabularx}
\normalsize

\clearpage
\onecolumn
\printbibliography[title={References}]
\end{document}
